% Complete reviewed Unit 5; editorial master: chapters/ch05/student.md
% Markdown SHA-256: d6b0fea4a1444547e3d857f24c2683455b372f9de908164a04125ac7b9416ded
% One-use initial assembly followed by complete content and page review.
\newcommand{\runningtitle}{Unit 5. Strong verbs IV–VII}
% Layout inherited from the reviewed previous course; set a running title per unit.
% This is a small hand-maintained TeX layout, not a conversion pipeline.
\RequirePackage{fix-cm}
\documentclass[11pt,twoside]{article}
\usepackage{fontspec}
\setmainfont{Cambria}[
  Renderer=HarfBuzz,
  BoldFont=cambriab.ttf,
  ItalicFont=cambriai.ttf,
  BoldItalicFont=cambriaz.ttf
]
\setmonofont{lmmono10-regular.otf}[BoldFont=lmmonolt10-bold.otf]
\usepackage[paperwidth=165mm,paperheight=241mm,inner=21mm,outer=19mm,
  top=22mm,bottom=20mm,headheight=14pt,headsep=7mm,footskip=10mm]{geometry}
\usepackage{xcolor,array,booktabs,longtable,tabularx}
\usepackage{enumitem,fancyhdr}
\usepackage[unicode,hidelinks]{hyperref}
\definecolor{ink}{HTML}{000000}
\definecolor{muted}{HTML}{4C5558}
\setlength{\parindent}{0pt}
\setlength{\parskip}{3.5pt}
\AtBeginDocument{\fontsize{11}{14.3}\selectfont}
\setlength{\tabcolsep}{4pt}
\renewcommand{\arraystretch}{1.10}
\setlist[itemize]{leftmargin=1.3em,label=\textbullet,itemsep=2pt,topsep=3pt,parsep=0pt}
\setlist[enumerate]{leftmargin=1.8em,itemsep=3pt,topsep=4pt,parsep=0pt}
\widowpenalty=10000
\clubpenalty=10000
\emergencystretch=1em
\raggedbottom
\pagestyle{fancy}
\fancyhf{}
\providecommand{\runningtitle}{Old English}
\fancyhead[LE,RO]{\small\itshape \runningtitle}
\fancyhead[LO,RE]{\small Old English}
\fancyfoot[LE,RO]{\small\thepage}
\renewcommand{\headrulewidth}{0.3pt}
\newcommand{\booktitle}[2]{
  {\fontsize{9}{11}\selectfont\color{muted}#1\par}
  \vspace{3pt}
  {\fontsize{20}{24}\selectfont\color{ink}#2\par}
  \vspace{6pt}
}
\newcommand{\source}[1]{\par{\footnotesize\color{muted}#1\par}}
\newcommand{\readingtext}[1]{
  \par\addvspace{5pt}
  \begingroup\leftskip=1.5em
    \fontsize{12}{16}\selectfont #1\par
  \endgroup\addvspace{5pt}
}
\makeatletter
\renewcommand{\section}{\@startsection{section}{1}{\z@}
  {-2.4ex \@plus -0.6ex \@minus -0.2ex}{0.8ex \@plus 0.2ex}
  {\normalfont\fontsize{12}{15}\selectfont\bfseries\color{ink}}}
\renewcommand{\subsection}{\@startsection{subsection}{2}{\z@}
  {-1.8ex \@plus -0.5ex \@minus -0.2ex}{0.5ex \@plus 0.2ex}
  {\normalfont\fontsize{11}{14}\selectfont\bfseries\color{ink}}}
\makeatother

\hyphenpenalty=10000
\exhyphenpenalty=10000
\emergencystretch=3em
\newcommand{\itempara}[2]{\par\addvspace{3pt}\noindent\hangindent=1.8em\hangafter=1\makebox[1.8em][l]{#1.}#2\par}
\newcommand{\reserve}[1]{\par\penalty-100\begingroup\dimen0=\pagegoal\advance\dimen0 by-\pagetotal\ifdim\dimen0<#1\newpage\fi\endgroup}
\begin{document}
\addfontfeatures{Ligatures=TeXOff}
\raggedright
\booktitle{UNIT 5}{Strong verbs IV–VII}

% mdparagraph:P001:begin
This unit completes the principal strong-verb classes with selected present and preterite indicative systems. Four further verbs from Classes I–III provide the forms needed for the complete fifth Cook exercise.
% mdparagraph:P001:end

% mdparagraph:P002:begin
In the tables, \textbf{1}, \textbf{2} and \textbf{3} mean first, second and third person; \textbf{sg.} and \textbf{pl.} mean singular and plural. \textbf{1–3 pl.} gives the same form for all three plural persons. Dual subjects also take plural verbs. Each verb’s paired finite tables cover all six present and all six preterite indicative cells. Preserve the printed quantities and the distinction between e and ę.
% mdparagraph:P002:end

\reserve{7\baselineskip}
% mdparagraph:P003:begin
\section*{1. Principal parts and endings}
% mdparagraph:P003:end

% mdparagraph:P004:begin
Strong verbs change their stem vowel in the preterite; the weak verbs learned earlier add a dental element. Four principal parts identify the strong stem series. The infinitive names the verb, the singular preterite gives the first/\allowbreak{}third singular stem, the plural preterite gives the plural stem, and the participle completes the cited series.
% mdparagraph:P004:end

% table: U05-ENDINGS
% mdtable:T01:begin
% course-table-style: 2026-10
% course-panel:T01-P01:begin
\par\begingroup
\fontsize{10.5}{13.2}\selectfont
\setlength{\LTleft}{0pt}\setlength{\LTright}{\fill}
\noindent
\begin{tabular}{@{}>{\raggedright\arraybackslash}p{39.459mm}>{\raggedright\arraybackslash}p{39.459mm}>{\raggedright\arraybackslash}p{39.459mm}@{}}
\toprule
\textbf{Person and number} & \textbf{Present} & \textbf{Preterite} \\
\midrule
1 sg. & -e, with the special contracted and -ie forms below & singular preterite stem; no added ending \\
2 sg. & -st or -est, with the listed stem changes & plural preterite stem + -e \\
3 sg. & -ð or -eð, with the listed stem changes & singular preterite stem; no added ending \\
1–3 pl. & -að, with the listed contracted and -iað forms & plural preterite stem + -on \\
\bottomrule
\end{tabular}
\par\endgroup
% course-panel:T01-P01:end
% mdtable:T01:end
% mdparagraph:P005:begin
The \textbf{second singular preterite uses the plural stem}. Thus beran has bær for the first and third singular, but bǣre beside plural bǣron. In other verbs the singular and plural stem vowels coincide; their different endings still identify the cells.
% mdparagraph:P005:end

% mdparagraph:P006:begin
The infinitive and participle columns are \textbf{supplied citations} throughout this unit. Consult them in the principal-part exercises. They do not require independent nonfinite recall or the use of infinitive, participial, perfect or passive constructions. A participle in a four-part entry is not another finite preterite.
% mdparagraph:P006:end

% mdparagraph:P007:begin
The second and third singular present often have a changed vowel and a reduced ending. Short e may appear as i; short a gives æ, short ea gives short ie, long ēa gives long īe, ā gives ǣ, and ō gives ē in the stated paradigms. These changes are lexical patterns, not an instruction to alter every similar-looking vowel. Dental combinations give forms such as -tst and -tt, and doubled stem consonants simplify before consonantal endings. The tables and notes give the resulting forms together.
% mdparagraph:P007:end

% teaching: NEW-U05-GRAM-PREFIXED-LEXEMES
% mdparagraph:P008:begin
\textbf{Gesēon, gescieppan and geswīcan} are learned as separate verbs. Their meanings and forms are stated explicitly. The prefix is not an instruction to create further verbs or to assign a fixed difference in meaning to every prefixed pair. The checked headword \textbf{gesēon} follows Cook’s dictionary ge-sēon, with the dictionary hyphen omitted; the preterite geseah keeps its printed spelling.
% mdparagraph:P008:end

\reserve{7\baselineskip}
% mdparagraph:P009:begin
\section*{2. Class IV}
% mdparagraph:P009:end

% teaching: NEW-U05-GRAM-CLASS-IV
% mdparagraph:P010:begin
The ordinary Class IV series is \textbf{e : æ : ǣ : o}. Beran shows it directly: the singular past has short æ, while the plural-based past has long ǣ. The fourth vowel belongs to the supplied participle.
% mdparagraph:P010:end

% teaching: NEW-U05-VOC-beran
% teaching: NEW-U05-CIT-beran
% teaching: NEW-U05-VOC-niman
% teaching: NEW-U05-CIT-niman
% teaching: NEW-U05-VOC-cuman
% teaching: NEW-U05-CIT-cuman
% table: U05-PARTS-IV
% mdtable:T02:begin
\par\begingroup
\fontsize{10.5}{13.2}\selectfont
\setlength{\LTleft}{0pt}\setlength{\LTright}{\fill}
\noindent
\begin{tabular}{@{}>{\raggedright\arraybackslash}p{23.14mm}>{\raggedright\arraybackslash}p{21.04mm}>{\raggedright\arraybackslash}p{22.72mm}>{\raggedright\arraybackslash}p{22.72mm}>{\raggedright\arraybackslash}p{23.14mm}@{}}
\toprule
\textbf{Infinitive¹} & \textbf{Meaning} & \textbf{Preterite sg.} & \textbf{Preterite pl.} & \textbf{Participle¹} \\
\midrule
beran & carry & bær & bǣron & boren \\
niman & take & nōm & nōmon & numen \\
cuman & come & cōm & cōmon & cumen \\
\bottomrule
\end{tabular}
\par\endgroup

% mdtable:T02:end
% mdparagraph:P011:begin
¹ Infinitive and participle entries are supplied citations. The middle columns identify the finite preterite stems; the complete person assignments follow.
% mdparagraph:P011:end

% teaching: NEW-U05-PAR-beran
% teaching: NEW-U05-PAR-niman
% teaching: NEW-U05-PAR-cuman
\reserve{8\baselineskip}
% mdparagraph:P012:begin
\textbf{Present indicative.}
% mdparagraph:P012:end

% table: U05-PRESENT-IV
% mdtable:T03:begin
% course-table-style: 2026-10
% course-panel:T03-P01:begin
\par\begingroup
\fontsize{10.5}{13.2}\selectfont
\setlength{\LTleft}{0pt}\setlength{\LTright}{\fill}
\noindent
\begin{tabular}{@{}>{\raggedright\arraybackslash}p{30.902mm}>{\raggedright\arraybackslash}p{28.221mm}>{\raggedright\arraybackslash}p{28.221mm}>{\raggedright\arraybackslash}p{28.221mm}@{}}
\toprule
\textbf{Person and number} & \textbf{beran} & \textbf{niman} & \textbf{cuman} \\
\midrule
1 sg. & bere & nime & cume \\
2 sg. & birst & nimst & cymst \\
3 sg. & birð & nimð & cymð \\
1–3 pl. & berað & nimað & cumað \\
\bottomrule
\end{tabular}
\par\endgroup
% course-panel:T03-P01:end
% mdtable:T03:end
\reserve{8\baselineskip}
% mdparagraph:P013:begin
\textbf{Preterite indicative.}
% mdparagraph:P013:end

% table: U05-PRETERITE-IV
% mdtable:T04:begin
% course-table-style: 2026-10
% course-panel:T04-P01:begin
\par\begingroup
\fontsize{10.5}{13.2}\selectfont
\setlength{\LTleft}{0pt}\setlength{\LTright}{\fill}
\noindent
\begin{tabular}{@{}>{\raggedright\arraybackslash}p{30.902mm}>{\raggedright\arraybackslash}p{28.221mm}>{\raggedright\arraybackslash}p{28.221mm}>{\raggedright\arraybackslash}p{28.221mm}@{}}
\toprule
\textbf{Person and number} & \textbf{beran} & \textbf{niman} & \textbf{cuman} \\
\midrule
1 sg. & bær & nōm & cōm \\
2 sg. & bǣre & nōme & cōme \\
3 sg. & bær & nōm & cōm \\
1–3 pl. & bǣron & nōmon & cōmon \\
\bottomrule
\end{tabular}
\par\endgroup
% course-panel:T04-P01:end
% mdtable:T04:end
% teaching: NEW-U05-GRAM-NASAL-IV
% mdparagraph:P014:begin
Learn \textbf{niman} and \textbf{cuman} separately within this class. This course selects Cook’s nōm, nōme, nōmon and cōm, cōme, cōmon. Cuman changes u to y in cymst and cymð; its first singular cume and plural cumað retain u. Niman keeps i in its present forms.
% mdparagraph:P014:end

% mdparagraph:P015:begin
\textbf{Alternatives.} Beran also has bierst and bierð. Wright records nam and nāmon beside the selected nōm and nōmon. Cook gives cwōm and cwōmon beside cōm and cōmon, and the participle citation cymen beside cumen. These alternatives are available for comparison; use the selected table when a task expressly requests it.
% mdparagraph:P015:end

\reserve{7\baselineskip}
% mdparagraph:P016:begin
\section*{3. Class V}
% mdparagraph:P016:end

% teaching: NEW-U05-GRAM-CLASS-V
% mdparagraph:P017:begin
The ordinary Class V series is \textbf{e : æ : ǣ : e}. It differs from Class IV in the participle vowel. Sprecan gives the ordinary series; several other verbs require their own present or preterite changes.
% mdparagraph:P017:end

% teaching: NEW-U05-VOC-sprecan
% teaching: NEW-U05-CIT-sprecan
% teaching: NEW-U05-VOC-cwethan
% teaching: NEW-U05-CIT-cwethan
% teaching: NEW-U05-VOC-giefan
% teaching: NEW-U05-CIT-giefan
% teaching: NEW-U05-VOC-biddan
% teaching: NEW-U05-CIT-biddan
% teaching: NEW-U05-VOC-licgan
% teaching: NEW-U05-CIT-licgan
% teaching: NEW-U05-VOC-sittan
% teaching: NEW-U05-CIT-sittan
% teaching: NEW-U05-VOC-etan
% teaching: NEW-U05-CIT-etan
% teaching: NEW-U05-VOC-seon
% teaching: NEW-U05-CIT-seon
% teaching: NEW-U05-VOC-geseon
% teaching: NEW-U05-CIT-geseon
% table: U05-PARTS-V
% mdtable:T05:begin
\par\begingroup
\fontsize{10.5}{13.2}\selectfont
\setlength{\LTleft}{0pt}\setlength{\LTright}{\fill}
\noindent
\begin{tabular}{@{}>{\raggedright\arraybackslash}p{22.88mm}>{\raggedright\arraybackslash}p{22.05mm}>{\raggedright\arraybackslash}p{22.47mm}>{\raggedright\arraybackslash}p{22.47mm}>{\raggedright\arraybackslash}p{22.88mm}@{}}
\toprule
\textbf{Infinitive¹} & \textbf{Meaning} & \textbf{Preterite sg.} & \textbf{Preterite pl.} & \textbf{Participle¹} \\
\midrule
sprecan & speak & spræc & sprǣcon & sprecen \\
cweðan & say & cwæð & cwǣdon & cweden \\
giefan & give & geaf & gēafon & giefen \\
biddan & ask; request & bæd & bǣdon & beden \\
licgan & lie; be at rest & læg & lǣgon & legen \\
sittan & sit & sæt & sǣton & seten \\
etan & eat & ǣt & ǣton & eten \\
sēon & see & seah & sāwon & sewen \\
gesēon & see & geseah & gesāwon & gesewen \\
\bottomrule
\end{tabular}
\par\endgroup

% mdtable:T05:end
% mdparagraph:P018:begin
¹ Infinitive and participle entries are supplied citations. The middle columns identify the finite preterite stems; the complete person assignments follow.
% mdparagraph:P018:end

% teaching: NEW-U05-PAR-sprecan
% teaching: NEW-U05-PAR-cwethan
% teaching: NEW-U05-PAR-giefan
% teaching: NEW-U05-PAR-biddan
% teaching: NEW-U05-PAR-licgan
% teaching: NEW-U05-PAR-sittan
% teaching: NEW-U05-PAR-etan
% teaching: NEW-U05-PAR-seon
% teaching: NEW-U05-PAR-geseon
\reserve{14\baselineskip}
% mdparagraph:P019:begin
\textbf{Present indicative.}
% mdparagraph:P019:end

% table: U05-PRESENT-V
% mdtable:T06:begin
% course-table-style: 2026-10
% course-panel:T06-P01:begin
\par\begingroup
\fontsize{10.5}{13.2}\selectfont
\setlength{\LTleft}{0pt}\setlength{\LTright}{\fill}
\noindent
\begin{tabular}{@{}>{\raggedright\arraybackslash}p{24.232mm}>{\raggedright\arraybackslash}p{22.13mm}>{\raggedright\arraybackslash}p{22.13mm}>{\raggedright\arraybackslash}p{22.13mm}>{\raggedright\arraybackslash}p{22.13mm}@{}}
\toprule
\textbf{Person and number} & \textbf{sprecan} & \textbf{cweðan} & \textbf{giefan} & \textbf{biddan} \\
\midrule
1 sg. & sprece & cweðe & giefe & bidde \\
2 sg. & spricst & cwist & giefst & bitst \\
3 sg. & spricð & cwið & giefð & bitt \\
1–3 pl. & sprecað & cweðað & giefað & biddað \\
\bottomrule
\end{tabular}
\par\endgroup
% course-panel:T06-P01:end
\par\addvspace{5pt}
% course-panel:T06-P02:begin
\par\begingroup
\fontsize{10.5}{13.2}\selectfont
\setlength{\LTleft}{0pt}\setlength{\LTright}{\fill}
\noindent
\begin{tabular}{@{}>{\raggedright\arraybackslash}p{24.232mm}>{\raggedright\arraybackslash}p{22.13mm}>{\raggedright\arraybackslash}p{22.13mm}>{\raggedright\arraybackslash}p{22.13mm}>{\raggedright\arraybackslash}p{22.13mm}@{}}
\toprule
\textbf{Person and number} & \textbf{licgan} & \textbf{sittan} & \textbf{etan} & \textbf{sēon} \\
\midrule
1 sg. & licge & sitte & ete & sēo \\
2 sg. & ligst & sitst & itst & siehst \\
3 sg. & ligð & sitt & itt & siehð \\
1–3 pl. & licgað & sittað & etað & sēoð \\
\bottomrule
\end{tabular}
\par\endgroup
% course-panel:T06-P02:end
\par\addvspace{5pt}
% course-panel:T06-P03:begin
\par\begingroup
\fontsize{10.5}{13.2}\selectfont
\setlength{\LTleft}{0pt}\setlength{\LTright}{\fill}
\noindent
\begin{tabular}{@{}>{\raggedright\arraybackslash}p{63.342mm}>{\raggedright\arraybackslash}p{57.846mm}@{}}
\toprule
\textbf{Person and number} & \textbf{gesēon} \\
\midrule
1 sg. & gesēo \\
2 sg. & gesiehst \\
3 sg. & gesiehð \\
1–3 pl. & gesēoð \\
\bottomrule
\end{tabular}
\par\endgroup
% course-panel:T06-P03:end
% mdtable:T06:end
\reserve{14\baselineskip}
% mdparagraph:P020:begin
\textbf{Preterite indicative.}
% mdparagraph:P020:end

% table: U05-PRETERITE-V
% mdtable:T07:begin
% course-table-style: 2026-10
% course-panel:T07-P01:begin
\par\begingroup
\fontsize{10.5}{13.2}\selectfont
\setlength{\LTleft}{0pt}\setlength{\LTright}{\fill}
\noindent
\begin{tabular}{@{}>{\raggedright\arraybackslash}p{24.232mm}>{\raggedright\arraybackslash}p{22.13mm}>{\raggedright\arraybackslash}p{22.13mm}>{\raggedright\arraybackslash}p{22.13mm}>{\raggedright\arraybackslash}p{22.13mm}@{}}
\toprule
\textbf{Person and number} & \textbf{sprecan} & \textbf{cweðan} & \textbf{giefan} & \textbf{biddan} \\
\midrule
1 sg. & spræc & cwæð & geaf & bæd \\
2 sg. & sprǣce & cwǣde & gēafe & bǣde \\
3 sg. & spræc & cwæð & geaf & bæd \\
1–3 pl. & sprǣcon & cwǣdon & gēafon & bǣdon \\
\bottomrule
\end{tabular}
\par\endgroup
% course-panel:T07-P01:end
\par\addvspace{5pt}
% course-panel:T07-P02:begin
\par\begingroup
\fontsize{10.5}{13.2}\selectfont
\setlength{\LTleft}{0pt}\setlength{\LTright}{\fill}
\noindent
\begin{tabular}{@{}>{\raggedright\arraybackslash}p{24.232mm}>{\raggedright\arraybackslash}p{22.13mm}>{\raggedright\arraybackslash}p{22.13mm}>{\raggedright\arraybackslash}p{22.13mm}>{\raggedright\arraybackslash}p{22.13mm}@{}}
\toprule
\textbf{Person and number} & \textbf{licgan} & \textbf{sittan} & \textbf{etan} & \textbf{sēon} \\
\midrule
1 sg. & læg & sæt & ǣt & seah \\
2 sg. & lǣge & sǣte & ǣte & sāwe \\
3 sg. & læg & sæt & ǣt & seah \\
1–3 pl. & lǣgon & sǣton & ǣton & sāwon \\
\bottomrule
\end{tabular}
\par\endgroup
% course-panel:T07-P02:end
\par\addvspace{5pt}
% course-panel:T07-P03:begin
\par\begingroup
\fontsize{10.5}{13.2}\selectfont
\setlength{\LTleft}{0pt}\setlength{\LTright}{\fill}
\noindent
\begin{tabular}{@{}>{\raggedright\arraybackslash}p{63.342mm}>{\raggedright\arraybackslash}p{57.846mm}@{}}
\toprule
\textbf{Person and number} & \textbf{gesēon} \\
\midrule
1 sg. & geseah \\
2 sg. & gesāwe \\
3 sg. & geseah \\
1–3 pl. & gesāwon \\
\bottomrule
\end{tabular}
\par\endgroup
% course-panel:T07-P03:end
% mdtable:T07:end
% teaching: NEW-U05-GRAM-CWETHAN-CONSONANTS
% mdparagraph:P021:begin
\textbf{Cweðan} has cwist and cwið in the selected second and third singular present. Its past singular cwæð retains ð, while cwǣde and cwǣdon have d; the supplied participle cweden also has d. Keep this consonant change with these stem positions.
% mdparagraph:P021:end

% teaching: NEW-U05-GRAM-PALATAL-V
% mdparagraph:P022:begin
\textbf{Giefan} has the palatal series giefan, geaf, gēafon, giefen. In particular, geaf has short ea, but gēafe and gēafon have long ēa. The present forms retain the selected ie spelling.
% mdparagraph:P022:end

% teaching: NEW-U05-GRAM-J-PRESENT-V
% mdparagraph:P023:begin
\textbf{Biddan, licgan and sittan} have doubled consonants in the present first singular and plural, but a single stem consonant before the second/\allowbreak{}third singular endings. The resulting forms are bitst/\allowbreak{}bitt, ligst/\allowbreak{}ligð and sitst/\allowbreak{}sitt. Their strong pasts use bæd/\allowbreak{}bǣdon, læg/\allowbreak{}lǣgon and sæt/\allowbreak{}sǣton. These are often called \textbf{j-presents}: the label explains a present-stem pattern, not a weak preterite.
% mdparagraph:P023:end

% teaching: NEW-U05-GRAM-ETAN-EXCEPTION
% mdparagraph:P024:begin
\textbf{Etan} has itst and itt in the present singular. The selected preterite ǣt has long ǣ, as do ǣte and ǣton. Its singular past therefore differs from the ordinary short-æ Class V pattern.
% mdparagraph:P024:end

% teaching: NEW-U05-GRAM-CONTRACT-V
% mdparagraph:P025:begin
\textbf{Sēon} and \textbf{gesēon} have contracted presents. Their first singular ends in -o and their plural in -oð. The second and third singular have \textbf{short ie}: siehst/\allowbreak{}siehð and gesiehst/\allowbreak{}gesiehð. Their plural-based past stems have āw: sāwe/\allowbreak{}sāwon and gesāwe/\allowbreak{}gesāwon. Preserve the difference between these forms and the singular seah/\allowbreak{}geseah.
% mdparagraph:P025:end

% mdparagraph:P026:begin
\textbf{Alternatives.} Third-singular bit, sit and it occur beside selected bitt, sitt and itt. Cook gives līst and līð beside ligst and ligð, and short æt beside selected ǣt; Wright also gives lāgon beside lǣgon. For sēon, Cook gives the participle segen beside sewen, and Wright records sǣgon and sawen beside sāwon and sewen. The corresponding prefixed citation alternative is gesegen beside gesewen. These spellings do not replace the selected tables or require extra nonfinite recall.
% mdparagraph:P026:end

\reserve{7\baselineskip}
% mdparagraph:P027:begin
\section*{4. Class VI}
% mdparagraph:P027:end

% teaching: NEW-U05-GRAM-CLASS-VI
% mdparagraph:P028:begin
The ordinary Class VI series is \textbf{a : ō : ō : a}. Faran shows it directly. The singular and plural past stems share long ō, and the second singular adds -e to that stem. Present mutation and the special present stems must still be learned separately.
% mdparagraph:P028:end

% teaching: NEW-U05-VOC-faran
% teaching: NEW-U05-CIT-faran
% teaching: NEW-U05-VOC-slean
% teaching: NEW-U05-CIT-slean
% teaching: NEW-U05-VOC-standan
% teaching: NEW-U05-CIT-standan
% teaching: NEW-U05-VOC-hebban
% teaching: NEW-U05-CIT-hebban
% teaching: NEW-U05-VOC-scieppan
% teaching: NEW-U05-CIT-scieppan
% teaching: NEW-U05-VOC-gescieppan
% teaching: NEW-U05-CIT-gescieppan
% teaching: NEW-U05-VOC-swerian
% teaching: NEW-U05-CIT-swerian
% table: U05-PARTS-VI
% mdtable:T08:begin
\par\begingroup
\fontsize{10.5}{13.2}\selectfont
\setlength{\LTleft}{0pt}\setlength{\LTright}{\fill}
\noindent
\begin{tabular}{@{}>{\raggedright\arraybackslash}p{23.14mm}>{\raggedright\arraybackslash}p{21.04mm}>{\raggedright\arraybackslash}p{22.72mm}>{\raggedright\arraybackslash}p{22.72mm}>{\raggedright\arraybackslash}p{23.14mm}@{}}
\toprule
\textbf{Infinitive¹} & \textbf{Meaning} & \textbf{Preterite sg.} & \textbf{Preterite pl.} & \textbf{Participle¹} \\
\midrule
faran & go; travel & fōr & fōron & faren \\
slēan & strike; smite & slōg & slōgon & slægen \\
standan & stand & stōd & stōdon & standen \\
hębban & raise; lift & hōf & hōfon & hafen \\
scieppan & create; shape & scēop & scēopon & sceapen \\
gescieppan & create & gescēop & gescēopon & gesceapen \\
swęrian & swear & swōr & swōron & sworen \\
\bottomrule
\end{tabular}
\par\endgroup

% mdtable:T08:end
% mdparagraph:P029:begin
¹ Infinitive and participle entries are supplied citations. The middle columns identify the finite preterite stems; the complete person assignments follow.
% mdparagraph:P029:end

% teaching: NEW-U05-PAR-faran
% teaching: NEW-U05-PAR-slean
% teaching: NEW-U05-PAR-standan
% teaching: NEW-U05-PAR-hebban
% teaching: NEW-U05-PAR-scieppan
% teaching: NEW-U05-PAR-gescieppan
% teaching: NEW-U05-PAR-swerian
\reserve{12\baselineskip}
% mdparagraph:P030:begin
\textbf{Present indicative.}
% mdparagraph:P030:end

% table: U05-PRESENT-VI
% mdtable:T09:begin
% course-table-style: 2026-10
% course-panel:T09-P01:begin
\par\begingroup
\fontsize{10.5}{13.2}\selectfont
\setlength{\LTleft}{0pt}\setlength{\LTright}{\fill}
\noindent
\begin{tabular}{@{}>{\raggedright\arraybackslash}p{24.232mm}>{\raggedright\arraybackslash}p{22.13mm}>{\raggedright\arraybackslash}p{22.13mm}>{\raggedright\arraybackslash}p{22.13mm}>{\raggedright\arraybackslash}p{22.13mm}@{}}
\toprule
\textbf{Person and number} & \textbf{faran} & \textbf{slēan} & \textbf{standan} & \textbf{hębban} \\
\midrule
1 sg. & fare & slēa & stande & hębbe \\
2 sg. & færst & sliehst & stęntst & hęfst \\
3 sg. & færð & sliehð & stęnt & hęfð \\
1–3 pl. & farað & slēað & standað & hębbað \\
\bottomrule
\end{tabular}
\par\endgroup
% course-panel:T09-P01:end
\par\addvspace{5pt}
% course-panel:T09-P02:begin
\par\begingroup
\fontsize{10.5}{13.2}\selectfont
\setlength{\LTleft}{0pt}\setlength{\LTright}{\fill}
\noindent
\begin{tabular}{@{}>{\raggedright\arraybackslash}p{30.902mm}>{\raggedright\arraybackslash}p{28.221mm}>{\raggedright\arraybackslash}p{28.221mm}>{\raggedright\arraybackslash}p{28.221mm}@{}}
\toprule
\textbf{Person and number} & \textbf{scieppan} & \textbf{gescieppan} & \textbf{swęrian} \\
\midrule
1 sg. & scieppe & gescieppe & swęrie \\
2 sg. & sciepst & gesciepst & swęrest \\
3 sg. & sciepð & gesciepð & swęreð \\
1–3 pl. & scieppað & gescieppað & swęriað \\
\bottomrule
\end{tabular}
\par\endgroup
% course-panel:T09-P02:end
% mdtable:T09:end
\reserve{12\baselineskip}
% mdparagraph:P031:begin
\textbf{Preterite indicative.}
% mdparagraph:P031:end

% table: U05-PRETERITE-VI
% mdtable:T10:begin
% course-table-style: 2026-10
% course-panel:T10-P01:begin
\par\begingroup
\fontsize{10.5}{13.2}\selectfont
\setlength{\LTleft}{0pt}\setlength{\LTright}{\fill}
\noindent
\begin{tabular}{@{}>{\raggedright\arraybackslash}p{24.232mm}>{\raggedright\arraybackslash}p{22.13mm}>{\raggedright\arraybackslash}p{22.13mm}>{\raggedright\arraybackslash}p{22.13mm}>{\raggedright\arraybackslash}p{22.13mm}@{}}
\toprule
\textbf{Person and number} & \textbf{faran} & \textbf{slēan} & \textbf{standan} & \textbf{hębban} \\
\midrule
1 sg. & fōr & slōg & stōd & hōf \\
2 sg. & fōre & slōge & stōde & hōfe \\
3 sg. & fōr & slōg & stōd & hōf \\
1–3 pl. & fōron & slōgon & stōdon & hōfon \\
\bottomrule
\end{tabular}
\par\endgroup
% course-panel:T10-P01:end
\par\addvspace{5pt}
% course-panel:T10-P02:begin
\par\begingroup
\fontsize{10.5}{13.2}\selectfont
\setlength{\LTleft}{0pt}\setlength{\LTright}{\fill}
\noindent
\begin{tabular}{@{}>{\raggedright\arraybackslash}p{30.902mm}>{\raggedright\arraybackslash}p{28.221mm}>{\raggedright\arraybackslash}p{28.221mm}>{\raggedright\arraybackslash}p{28.221mm}@{}}
\toprule
\textbf{Person and number} & \textbf{scieppan} & \textbf{gescieppan} & \textbf{swęrian} \\
\midrule
1 sg. & scēop & gescēop & swōr \\
2 sg. & scēope & gescēope & swōre \\
3 sg. & scēop & gescēop & swōr \\
1–3 pl. & scēopon & gescēopon & swōron \\
\bottomrule
\end{tabular}
\par\endgroup
% course-panel:T10-P02:end
% mdtable:T10:end
% teaching: NEW-U05-GRAM-CONTRACT-VI
% mdparagraph:P032:begin
\textbf{Slēan} has a contracted present: slēa, sliehst, sliehð, slēað. The ie in sliehst and sliehð is short. The selected past has g throughout: slōg, slōge, slōgon. The supplied participle is slægen.
% mdparagraph:P032:end

% teaching: NEW-U05-GRAM-STANDAN-NASAL
% mdparagraph:P033:begin
\textbf{Standan} keeps n in the present but loses it in the past stōd/\allowbreak{}stōdon. Its changed present singulars are stęntst and stęnt. This is the ordinary verb of standing; the earlier reading’s opposing-winds construction requires its own supplied explanation.
% mdparagraph:P033:end

% teaching: NEW-U05-GRAM-J-PRESENT-VI
% mdparagraph:P034:begin
\textbf{Hębban} and \textbf{scieppan} have doubled present consonants in hębbe/\allowbreak{}hębbað and scieppe/\allowbreak{}scieppað. The selected second and third singulars have f in hęfst/\allowbreak{}hęfð and single p in sciepst/\allowbreak{}sciepð. Their strong pasts are hōf/\allowbreak{}hōfon and scēop/\allowbreak{}scēopon, with supplied participles hafen and sceapen. Gescieppan follows its separately given prefixed table.
% mdparagraph:P034:end

% mdparagraph:P035:begin
\textbf{Swęrian} has swęrie, swęrest, swęreð, swęriað, retaining full endings in the second and third singular. Its past is nevertheless strong: swōr, swōre, swōron. An infinitive in -ian does not by itself identify a weak verb. Keep the full -est and -eð in this selected present.
% mdparagraph:P035:end

% mdparagraph:P036:begin
\textbf{Alternatives.} Cook gives scōp/\allowbreak{}scōpon alongside scēop/\allowbreak{}scēopon. The corresponding prefix-preserving alternatives are gescōp/\allowbreak{}gescōpon. Cook also gives the participle slegen beside slægen; Wright records slōh beside singular slōg and the participle slagen. The selected cells remain those in the tables.
% mdparagraph:P036:end

\reserve{7\baselineskip}
% mdparagraph:P037:begin
\section*{5. Class VII}
% mdparagraph:P037:end

% mdparagraph:P038:begin
Cook groups these verbs as \textbf{reduplicating verbs}; they are commonly called Class VII. For present purposes, learn two past-vowel groups, ēo and ē. Within each verb the selected preterite stem is the same in singular and plural. The present and participle usually share their vowel, but the complete entries preserve exceptions such as wēpan/\allowbreak{}wōpen and fōn/\allowbreak{}fangen.
% mdparagraph:P038:end

\reserve{7\baselineskip}
% mdparagraph:P039:begin
\subsection*{Preterites with ēo}
% mdparagraph:P039:end

% teaching: NEW-U05-GRAM-CLASS-VII-EO
% teaching: NEW-U05-VOC-feallan
% teaching: NEW-U05-CIT-feallan
% teaching: NEW-U05-VOC-healdan
% teaching: NEW-U05-CIT-healdan
% teaching: NEW-U05-VOC-weaxan
% teaching: NEW-U05-CIT-weaxan
% teaching: NEW-U05-VOC-cnawan
% teaching: NEW-U05-CIT-cnawan
% teaching: NEW-U05-VOC-heawan
% teaching: NEW-U05-CIT-heawan
% teaching: NEW-U05-VOC-flowan
% teaching: NEW-U05-CIT-flowan
% teaching: NEW-U05-VOC-wepan
% teaching: NEW-U05-CIT-wepan
% table: U05-PARTS-VII-EO
% mdtable:T11:begin
\par\begingroup
\fontsize{10.5}{13.2}\selectfont
\setlength{\LTleft}{0pt}\setlength{\LTright}{\fill}
\noindent
\begin{tabular}{@{}>{\raggedright\arraybackslash}p{23.14mm}>{\raggedright\arraybackslash}p{21.04mm}>{\raggedright\arraybackslash}p{22.72mm}>{\raggedright\arraybackslash}p{22.72mm}>{\raggedright\arraybackslash}p{23.14mm}@{}}
\toprule
\textbf{Infinitive¹} & \textbf{Meaning} & \textbf{Preterite sg.} & \textbf{Preterite pl.} & \textbf{Participle¹} \\
\midrule
feallan & fall & fēoll & fēollon & feallen \\
healdan & hold & hēold & hēoldon & healden \\
weaxan & grow & wēox & wēoxon & weaxen \\
cnāwan & know & cnēow & cnēowon & cnāwen \\
hēawan & hew & hēow & hēowon & hēawen \\
flōwan & flow & flēow & flēowon & flōwen \\
wēpan & weep & wēop & wēopon & wōpen \\
\bottomrule
\end{tabular}
\par\endgroup

% mdtable:T11:end
% mdparagraph:P040:begin
¹ Infinitive and participle entries are supplied citations. The middle columns identify the finite preterite stems; the complete person assignments follow.
% mdparagraph:P040:end

% teaching: NEW-U05-PAR-feallan
% teaching: NEW-U05-PAR-healdan
% teaching: NEW-U05-PAR-weaxan
% teaching: NEW-U05-PAR-cnawan
% teaching: NEW-U05-PAR-heawan
% teaching: NEW-U05-PAR-flowan
% teaching: NEW-U05-PAR-wepan
\reserve{12\baselineskip}
% mdparagraph:P041:begin
\textbf{Present indicative.}
% mdparagraph:P041:end

% table: U05-PRESENT-VII-EO
% mdtable:T12:begin
% course-table-style: 2026-10
% course-panel:T12-P01:begin
\par\begingroup
\fontsize{10.5}{13.2}\selectfont
\setlength{\LTleft}{0pt}\setlength{\LTright}{\fill}
\noindent
\begin{tabular}{@{}>{\raggedright\arraybackslash}p{24.232mm}>{\raggedright\arraybackslash}p{22.13mm}>{\raggedright\arraybackslash}p{22.13mm}>{\raggedright\arraybackslash}p{22.13mm}>{\raggedright\arraybackslash}p{22.13mm}@{}}
\toprule
\textbf{Person and number} & \textbf{feallan} & \textbf{healdan} & \textbf{weaxan} & \textbf{cnāwan} \\
\midrule
1 sg. & fealle & healde & weaxe & cnāwe \\
2 sg. & fielst & hieltst & wiext & cnǣwst \\
3 sg. & fielð & hielt & wiext & cnǣwð \\
1–3 pl. & feallað & healdað & weaxað & cnāwað \\
\bottomrule
\end{tabular}
\par\endgroup
% course-panel:T12-P01:end
\par\addvspace{5pt}
% course-panel:T12-P02:begin
\par\begingroup
\fontsize{10.5}{13.2}\selectfont
\setlength{\LTleft}{0pt}\setlength{\LTright}{\fill}
\noindent
\begin{tabular}{@{}>{\raggedright\arraybackslash}p{30.902mm}>{\raggedright\arraybackslash}p{28.221mm}>{\raggedright\arraybackslash}p{28.221mm}>{\raggedright\arraybackslash}p{28.221mm}@{}}
\toprule
\textbf{Person and number} & \textbf{hēawan} & \textbf{flōwan} & \textbf{wēpan} \\
\midrule
1 sg. & hēawe & flōwe & wēpe \\
2 sg. & hīewst & flēwst & wēpst \\
3 sg. & hīewð & flēwð & wēpð \\
1–3 pl. & hēawað & flōwað & wēpað \\
\bottomrule
\end{tabular}
\par\endgroup
% course-panel:T12-P02:end
% mdtable:T12:end
\reserve{12\baselineskip}
% mdparagraph:P042:begin
\textbf{Preterite indicative.}
% mdparagraph:P042:end

% table: U05-PRETERITE-VII-EO
% mdtable:T13:begin
% course-table-style: 2026-10
% course-panel:T13-P01:begin
\par\begingroup
\fontsize{10.5}{13.2}\selectfont
\setlength{\LTleft}{0pt}\setlength{\LTright}{\fill}
\noindent
\begin{tabular}{@{}>{\raggedright\arraybackslash}p{24.232mm}>{\raggedright\arraybackslash}p{22.13mm}>{\raggedright\arraybackslash}p{22.13mm}>{\raggedright\arraybackslash}p{22.13mm}>{\raggedright\arraybackslash}p{22.13mm}@{}}
\toprule
\textbf{Person and number} & \textbf{feallan} & \textbf{healdan} & \textbf{weaxan} & \textbf{cnāwan} \\
\midrule
1 sg. & fēoll & hēold & wēox & cnēow \\
2 sg. & fēolle & hēolde & wēoxe & cnēowe \\
3 sg. & fēoll & hēold & wēox & cnēow \\
1–3 pl. & fēollon & hēoldon & wēoxon & cnēowon \\
\bottomrule
\end{tabular}
\par\endgroup
% course-panel:T13-P01:end
\par\addvspace{5pt}
% course-panel:T13-P02:begin
\par\begingroup
\fontsize{10.5}{13.2}\selectfont
\setlength{\LTleft}{0pt}\setlength{\LTright}{\fill}
\noindent
\begin{tabular}{@{}>{\raggedright\arraybackslash}p{30.902mm}>{\raggedright\arraybackslash}p{28.221mm}>{\raggedright\arraybackslash}p{28.221mm}>{\raggedright\arraybackslash}p{28.221mm}@{}}
\toprule
\textbf{Person and number} & \textbf{hēawan} & \textbf{flōwan} & \textbf{wēpan} \\
\midrule
1 sg. & hēow & flēow & wēop \\
2 sg. & hēowe & flēowe & wēope \\
3 sg. & hēow & flēow & wēop \\
1–3 pl. & hēowon & flēowon & wēopon \\
\bottomrule
\end{tabular}
\par\endgroup
% course-panel:T13-P02:end
% mdtable:T13:end
% teaching: NEW-U05-GRAM-VII-PRESENT-CHANGES
% mdparagraph:P043:begin
In the present second and third singular, \textbf{feallan, healdan and weaxan} change short ea to short ie. Healdan also has the dental results hieltst and hielt. Stem-final x combines with either singular ending as xt: \textbf{wiext is both second and third singular} of weaxan. Feallan reduces ll before its endings, giving fielst and fielð.
% mdparagraph:P043:end

% mdparagraph:P044:begin
Cnāwan changes ā to ǣ in cnǣwst/\allowbreak{}cnǣwð; hēawan changes long ēa to long īe in hīewst/\allowbreak{}hīewð; flōwan changes ō to ē in flēwst/\allowbreak{}flēwð. Wēpan keeps ē in wēpst/\allowbreak{}wēpð. In the preterite, cnēow/\allowbreak{}cnēowon and hēow/\allowbreak{}hēowon retain their printed consonants and quantities; do not transfer a consonant from another verb’s past.
% mdparagraph:P044:end

\reserve{7\baselineskip}
% mdparagraph:P045:begin
\subsection*{Preterites with ē}
% mdparagraph:P045:end

% teaching: NEW-U05-GRAM-CLASS-VII-E
% teaching: NEW-U05-VOC-laetan
% teaching: NEW-U05-CIT-laetan
% teaching: NEW-U05-VOC-hatan
% teaching: NEW-U05-CIT-hatan
% teaching: NEW-U05-VOC-fon
% teaching: NEW-U05-CIT-fon
% table: U05-PARTS-VII-E
% mdtable:T14:begin
\par\begingroup
\fontsize{10.5}{13.2}\selectfont
\setlength{\LTleft}{0pt}\setlength{\LTright}{\fill}
\noindent
\begin{tabular}{@{}>{\raggedright\arraybackslash}p{23.14mm}>{\raggedright\arraybackslash}p{21.04mm}>{\raggedright\arraybackslash}p{22.72mm}>{\raggedright\arraybackslash}p{22.72mm}>{\raggedright\arraybackslash}p{23.14mm}@{}}
\toprule
\textbf{Infinitive¹} & \textbf{Meaning} & \textbf{Preterite sg.} & \textbf{Preterite pl.} & \textbf{Participle¹} \\
\midrule
lǣtan & let & lēt & lēton & lǣten \\
hātan & call; name & hēt & hēton & hāten \\
fōn & seize; catch & fēng & fēngon & fangen \\
\bottomrule
\end{tabular}
\par\endgroup

% mdtable:T14:end
% mdparagraph:P046:begin
¹ Infinitive and participle entries are supplied citations. The middle columns identify the finite preterite stems; the complete person assignments follow.
% mdparagraph:P046:end

% teaching: NEW-U05-PAR-laetan
% teaching: NEW-U05-PAR-hatan
% teaching: NEW-U05-PAR-fon
\reserve{8\baselineskip}
% mdparagraph:P047:begin
\textbf{Present indicative.}
% mdparagraph:P047:end

% table: U05-PRESENT-VII-E
% mdtable:T15:begin
% course-table-style: 2026-10
% course-panel:T15-P01:begin
\par\begingroup
\fontsize{10.5}{13.2}\selectfont
\setlength{\LTleft}{0pt}\setlength{\LTright}{\fill}
\noindent
\begin{tabular}{@{}>{\raggedright\arraybackslash}p{30.902mm}>{\raggedright\arraybackslash}p{28.221mm}>{\raggedright\arraybackslash}p{28.221mm}>{\raggedright\arraybackslash}p{28.221mm}@{}}
\toprule
\textbf{Person and number} & \textbf{lǣtan} & \textbf{hātan} & \textbf{fōn} \\
\midrule
1 sg. & lǣte & hāte & fō \\
2 sg. & lǣtst & hǣtst & fēhst \\
3 sg. & lǣtt & hǣtt & fēhð \\
1–3 pl. & lǣtað & hātað & fōð \\
\bottomrule
\end{tabular}
\par\endgroup
% course-panel:T15-P01:end
% mdtable:T15:end
\reserve{8\baselineskip}
% mdparagraph:P048:begin
\textbf{Preterite indicative.}
% mdparagraph:P048:end

% table: U05-PRETERITE-VII-E
% mdtable:T16:begin
% course-table-style: 2026-10
% course-panel:T16-P01:begin
\par\begingroup
\fontsize{10.5}{13.2}\selectfont
\setlength{\LTleft}{0pt}\setlength{\LTright}{\fill}
\noindent
\begin{tabular}{@{}>{\raggedright\arraybackslash}p{30.902mm}>{\raggedright\arraybackslash}p{28.221mm}>{\raggedright\arraybackslash}p{28.221mm}>{\raggedright\arraybackslash}p{28.221mm}@{}}
\toprule
\textbf{Person and number} & \textbf{lǣtan} & \textbf{hātan} & \textbf{fōn} \\
\midrule
1 sg. & lēt & hēt & fēng \\
2 sg. & lēte & hēte & fēnge \\
3 sg. & lēt & hēt & fēng \\
1–3 pl. & lēton & hēton & fēngon \\
\bottomrule
\end{tabular}
\par\endgroup
% course-panel:T16-P01:end
% mdtable:T16:end
% mdparagraph:P049:begin
\textbf{Lǣtan} has lǣtst and lǣtt in the present; \textbf{hātan} changes ā to ǣ in hǣtst and hǣtt. Their preterites are lēt/\allowbreak{}lēton and hēt/\allowbreak{}hēton. Third-singular lǣt and hǣt are documented alternatives. Wright also records the older preterite hēht; it is not an additional recall requirement here.
% mdparagraph:P049:end

% teaching: NEW-U05-GRAM-CONTRACT-VII
% mdparagraph:P050:begin
\textbf{Fōn} has the contracted present fō, fēhst, fēhð, fōð and the past fēng, fēnge, fēngon. The nasal consonant belongs to the past stem. The participle fangen remains a supplied citation.
% mdparagraph:P050:end

% mdparagraph:P051:begin
The hātan table is active indicative. The special passive form hātte, commands with infinitives, and the reading’s construction with an object and a name are separate matters. The naming phrases in §§11–12 remain supplied in full.
% mdparagraph:P051:end

\reserve{7\baselineskip}
% mdparagraph:P052:begin
\section*{6. Four additions to Classes I–III}
% mdparagraph:P052:end

% teaching: NEW-U05-GRAM-BRIDGES-I-III
% mdparagraph:P053:begin
The following verbs use classes already learned, but their complete lexical systems are new here. Earlier supplied occurrences do not replace learning these tables. Geswīcan and scīnan follow Class I, flēogan Class II, and swimman the nasal pattern of Class III.
% mdparagraph:P053:end

% teaching: NEW-U05-VOC-geswican
% teaching: NEW-U05-CIT-geswican
% teaching: NEW-U05-VOC-scinan
% teaching: NEW-U05-CIT-scinan
% teaching: NEW-U05-VOC-fleogan
% teaching: NEW-U05-CIT-fleogan
% teaching: NEW-U05-VOC-swimman
% teaching: NEW-U05-CIT-swimman
% table: U05-PARTS-I-III
% mdtable:T17:begin
\par\begingroup
\fontsize{10.5}{13.2}\selectfont
\setlength{\LTleft}{0pt}\setlength{\LTright}{\fill}
\noindent
\begin{tabular}{@{}>{\raggedright\arraybackslash}p{23.14mm}>{\raggedright\arraybackslash}p{21.04mm}>{\raggedright\arraybackslash}p{22.72mm}>{\raggedright\arraybackslash}p{22.72mm}>{\raggedright\arraybackslash}p{23.14mm}@{}}
\toprule
\textbf{Infinitive¹} & \textbf{Meaning} & \textbf{Preterite sg.} & \textbf{Preterite pl.} & \textbf{Participle¹} \\
\midrule
geswīcan & cease & geswāc & geswicon & geswicen \\
scīnan & shine & scān & scinon & scinen \\
flēogan & fly & flēag & flugon & flogen \\
swimman & swim & swamm & swummon & swummen \\
\bottomrule
\end{tabular}
\par\endgroup

% mdtable:T17:end
% mdparagraph:P054:begin
¹ Infinitive and participle entries are supplied citations. The middle columns identify the finite preterite stems; the complete person assignments follow.
% mdparagraph:P054:end

% teaching: NEW-U05-PAR-geswican
% teaching: NEW-U05-PAR-scinan
% teaching: NEW-U05-PAR-fleogan
% teaching: NEW-U05-PAR-swimman
\reserve{9\baselineskip}
% mdparagraph:P055:begin
\textbf{Present indicative.}
% mdparagraph:P055:end

% table: U05-PRESENT-I-III
% mdtable:T18:begin
% course-table-style: 2026-10
% course-panel:T18-P01:begin
\par\begingroup
\fontsize{10.5}{13.2}\selectfont
\setlength{\LTleft}{0pt}\setlength{\LTright}{\fill}
\noindent
\begin{tabular}{@{}>{\raggedright\arraybackslash}p{24.232mm}>{\raggedright\arraybackslash}p{22.13mm}>{\raggedright\arraybackslash}p{22.13mm}>{\raggedright\arraybackslash}p{22.13mm}>{\raggedright\arraybackslash}p{22.13mm}@{}}
\toprule
\textbf{Person and number} & \textbf{geswīcan} & \textbf{scīnan} & \textbf{flēogan} & \textbf{swimman} \\
\midrule
1 sg. & geswīce & scīne & flēoge & swimme \\
2 sg. & geswīcst & scīnst & flīehst & swimst \\
3 sg. & geswīcð & scīnð & flīehð & swimð \\
1–3 pl. & geswīcað & scīnað & flēogað & swimmað \\
\bottomrule
\end{tabular}
\par\endgroup
% course-panel:T18-P01:end
% mdtable:T18:end
\reserve{9\baselineskip}
% mdparagraph:P056:begin
\textbf{Preterite indicative.}
% mdparagraph:P056:end

% table: U05-PRETERITE-I-III
% mdtable:T19:begin
% course-table-style: 2026-10
% course-panel:T19-P01:begin
\par\begingroup
\fontsize{10.5}{13.2}\selectfont
\setlength{\LTleft}{0pt}\setlength{\LTright}{\fill}
\noindent
\begin{tabular}{@{}>{\raggedright\arraybackslash}p{24.232mm}>{\raggedright\arraybackslash}p{22.13mm}>{\raggedright\arraybackslash}p{22.13mm}>{\raggedright\arraybackslash}p{22.13mm}>{\raggedright\arraybackslash}p{22.13mm}@{}}
\toprule
\textbf{Person and number} & \textbf{geswīcan} & \textbf{scīnan} & \textbf{flēogan} & \textbf{swimman} \\
\midrule
1 sg. & geswāc & scān & flēag & swamm \\
2 sg. & geswice & scine & fluge & swumme \\
3 sg. & geswāc & scān & flēag & swamm \\
1–3 pl. & geswicon & scinon & flugon & swummon \\
\bottomrule
\end{tabular}
\par\endgroup
% course-panel:T19-P01:end
% mdtable:T19:end
% mdparagraph:P057:begin
Geswīcan and scīnan have long ī in the present, long ā in the first/\allowbreak{}third singular past, and short i in the plural-based past. Flēogan has ēo in the first singular and plural present, long īe in flīehst/\allowbreak{}flīehð, and u in fluge/\allowbreak{}flugon. Swimman has short a in swamm but u in swumme/\allowbreak{}swummon; mm reduces before its present consonantal endings.
% mdparagraph:P057:end

% mdparagraph:P058:begin
\textbf{Alternatives.} Cook records flīegst and flīegð beside selected flīehst and flīehð. The dictionary also attests swimman past alternatives swom and swommon. The tables select the a/\allowbreak{}u series for swimman and the h spellings for flēogan.
% mdparagraph:P058:end

\reserve{7\baselineskip}
% mdparagraph:P059:begin
\section*{7. Ordinary uses}
% mdparagraph:P059:end

% mdparagraph:P060:begin
The glosses in the principal-part tables identify the verbs. The constructions below supply the ordinary uses needed in sentences. For other verbs, the practice asks for labelled person, number and tense forms; a gloss alone does not specify which objects or clauses may follow.
% mdparagraph:P060:end

% teaching: NEW-U05-GRAM-CARRY-TAKE-OBJECTS
% mdparagraph:P061:begin
\textbf{Beran}, “carry,” and \textbf{niman}, “take,” take an accusative person or thing. The subject controls the finite verb. The following eight examples are new teaching material.
% mdparagraph:P061:end

% example: NEW-U05-E01
% mdparagraph:P062:begin
Ðū birst ðone męte.
% mdparagraph:P062:end

% mdparagraph:P063:begin
You (singular) carry the food.
% mdparagraph:P063:end

% example: NEW-U05-E02
% mdparagraph:P064:begin
Hē nimð ðone fisc.
% mdparagraph:P064:end

% mdparagraph:P065:begin
He takes the fish.
% mdparagraph:P065:end

% teaching: NEW-U05-GRAM-GIVING-RECIPIENT
% mdparagraph:P066:begin
With \textbf{giefan}, the thing given is accusative and the recipient is dative. English may express the recipient with “to”; Old English expresses the relation here by the case of the recipient itself.
% mdparagraph:P066:end

% example: NEW-U05-E03
% mdparagraph:P067:begin
Ic giefe ðǣm męnn ðone męte.
% mdparagraph:P067:end

% mdparagraph:P068:begin
I give the man the food.
% mdparagraph:P068:end

% teaching: NEW-U05-GRAM-INTRANSITIVE-POSTURE
% mdparagraph:P069:begin
\textbf{Sittan}, “sit,” \textbf{licgan}, “lie; be at rest,” and \textbf{standan}, “stand,” describe the subject’s posture. A locative \textbf{on + dative} phrase can state where that subject is. No object is needed for these ordinary posture uses.
% mdparagraph:P069:end

% example: NEW-U05-E04
% mdparagraph:P070:begin
Se mann ligð on ðǣm hofe.
% mdparagraph:P070:end

% mdparagraph:P071:begin
The man lies in the dwelling.
% mdparagraph:P071:end

% example: NEW-U05-E05
% mdparagraph:P072:begin
Se mann stęnt on ðǣm strande.
% mdparagraph:P072:end

% mdparagraph:P073:begin
The man stands on the shore.
% mdparagraph:P073:end

% teaching: NEW-U05-GRAM-INTRANSITIVE-MOTION
% teaching: NEW-U05-GRAM-BRIDGE-INTRANSITIVE
% mdparagraph:P074:begin
\textbf{Feallan}, “fall,” \textbf{weaxan}, “grow,” \textbf{flēogan}, “fly,” \textbf{scīnan}, “shine,” and \textbf{swimman}, “swim,” have ordinary intransitive uses. The growing tree, flying bird, shining star or swimming fish is the subject. A locative on-phrase may state where the swimming occurs.
% mdparagraph:P074:end

% example: NEW-U05-E06
% mdparagraph:P075:begin
Ðæt trēo fielð.
% mdparagraph:P075:end

% mdparagraph:P076:begin
The tree falls.
% mdparagraph:P076:end

% teaching: NEW-U05-GRAM-EAT
% teaching: NEW-U05-GRAM-SEE-OBJECT
% teaching: NEW-U05-GRAM-STRIKE-HOLD-OBJECTS
% teaching: NEW-U05-GRAM-U05-CREATE-OBJECT
% mdparagraph:P077:begin
\textbf{Etan}, “eat,” takes an accusative food object; it also occurs with the food understood. \textbf{Sēon} and \textbf{gesēon}, “see,” take an accusative person or thing in the use practised here. \textbf{Slēan}, “strike; smite,” and \textbf{healdan}, “hold,” likewise take an accusative object. “Strike” does not by itself say that the object is killed. \textbf{Scieppan} and \textbf{gescieppan} take the created being or thing as accusative object in the Creation sentences.
% mdparagraph:P077:end

% teaching: NEW-U05-GRAM-CESSATION-GENITIVE
% mdparagraph:P078:begin
\textbf{Geswīcan}, “cease,” takes the \textbf{genitive} of the work or action from which one ceases. Keep this government when English uses “cease from.” An agreeing possessive takes the case, number and gender of the work noun; the personal genitives his and hiera remain the possessor forms already learned.
% mdparagraph:P078:end

% example: NEW-U05-E07
% mdparagraph:P079:begin
Wē geswīcað ūres weorces.
% mdparagraph:P079:end

% mdparagraph:P080:begin
We cease from our work.
% mdparagraph:P080:end

% teaching: NEW-U05-GRAM-DIRECT-QUOTATION
% mdparagraph:P081:begin
\textbf{Cweðan}, “say,” can introduce direct quotation. The words inside the quotation form their own clause; a first-person pronoun there denotes the quoted speaker. Content clauses introduced by “that” remain separately supplied in the reading.
% mdparagraph:P081:end

% example: NEW-U05-E08
% mdparagraph:P082:begin
Se mann cwið: “Ic sēo ðone fugol.”
% mdparagraph:P082:end

% mdparagraph:P083:begin
The man says, “I see the bird.”
% mdparagraph:P083:end

% teaching: NEW-U05-GRAM-SIMPLE-NEGATION
% mdparagraph:P084:begin
For these ordinary indicative clauses, put \textbf{ne before the finite verb} to negate the statement. A past negative still uses the appropriate strong preterite; English “did not” does not add a separate Old English auxiliary here. Subjects in the dual take plural verb forms, while dual possessives still identify two possessors and agree with the possessed noun.
% mdparagraph:P084:end

% mdparagraph:P085:begin
No new productive noun pattern is needed. Continue the complete noun, adjective, article, pronoun and possessive paradigms of Units 1–4. In particular, the learned masculine heofon supplies plural “heavens”; the reading’s weak feminine heofone receives separate local help.
% mdparagraph:P085:end

\reserve{7\baselineskip}
% mdparagraph:P086:begin
\section*{8. Form and clause practice}
% mdparagraph:P086:end

% mdparagraph:P087:begin
P, F and C are new teaching exercises. \textbf{Selected} means the form in this unit’s main tables. Consult the printed principal parts where the prompt supplies them; the nonfinite entries are never independent recall targets. In translation tasks, retain each stated person, number, tense and case relation.
% mdparagraph:P087:end

\reserve{16\baselineskip}
% mdparagraph:P088:begin
\subsection*{Principal parts — P1–P4}
% mdparagraph:P088:end

\reserve{5\baselineskip}
\par\begin{minipage}{\linewidth}\raggedright
% mdparagraph:P089:begin
\textbf{P1.}
% mdparagraph:P089:end

% item: NEW-U05-P01
% mdparagraph:P090:begin
The printed entry gives beran, bær, bǣron, boren. Label the four principal-part positions. Which entries are finite, and which are nonfinite citations?
% mdparagraph:P090:end

\end{minipage}\par\addvspace{4pt}

\par\begin{minipage}{\linewidth}\raggedright
% mdparagraph:P091:begin
\textbf{P2.}
% mdparagraph:P091:end

% item: NEW-U05-P02
% mdparagraph:P092:begin
Compare geaf and sealde, both meaning gave. Identify the strong form and the weak form; explain the difference in preterite formation.
% mdparagraph:P092:end

\end{minipage}\par\addvspace{4pt}

\par\begin{minipage}{\linewidth}\raggedright
% mdparagraph:P093:begin
\textbf{P3.}
% mdparagraph:P093:end

% item: NEW-U05-P03
% mdparagraph:P094:begin
For beran, give the preterite indicative forms for I, you (singular), he, and we. Identify the forms using the preterite plural stem.
% mdparagraph:P094:end

\end{minipage}\par\addvspace{4pt}

\par\begin{minipage}{\linewidth}\raggedright
% mdparagraph:P095:begin
\textbf{P4.}
% mdparagraph:P095:end

% item: NEW-U05-P04
% mdparagraph:P096:begin
The printed entry gives gescieppan, gescēop, gescēopon, gesceapen. Using it, give you (singular) created and we created. Is gesceapen itself a finite preterite?
% mdparagraph:P096:end

\end{minipage}\par\addvspace{4pt}

\reserve{7\baselineskip}
% mdparagraph:P097:begin
\subsection*{Finite forms — F1–F24}
% mdparagraph:P097:end

\reserve{5\baselineskip}
\par\begin{minipage}{\linewidth}\raggedright
% mdparagraph:P098:begin
\textbf{F1.}
% mdparagraph:P098:end

% item: NEW-U05-F01
% mdparagraph:P099:begin
Give the selected indicative forms of beran: present 2sg; preterite 1sg and 2sg.
% mdparagraph:P099:end

\end{minipage}\par\addvspace{4pt}

\par\begin{minipage}{\linewidth}\raggedright
% mdparagraph:P100:begin
\textbf{F2.}
% mdparagraph:P100:end

% item: NEW-U05-F02
% mdparagraph:P101:begin
Give the selected indicative forms of niman: present 3sg; preterite 3sg and 1pl.
% mdparagraph:P101:end

\end{minipage}\par\addvspace{4pt}

\par\begin{minipage}{\linewidth}\raggedright
% mdparagraph:P102:begin
\textbf{F3.}
% mdparagraph:P102:end

% item: NEW-U05-F03
% mdparagraph:P103:begin
Give the selected indicative forms of cuman: present 1sg, 2sg and 3sg; preterite 3pl.
% mdparagraph:P103:end

\end{minipage}\par\addvspace{4pt}

\par\begin{minipage}{\linewidth}\raggedright
% mdparagraph:P104:begin
\textbf{F4.}
% mdparagraph:P104:end

% item: NEW-U05-F04
% mdparagraph:P105:begin
Give the selected indicative forms of sprecan: present 2sg and 3sg; preterite 2sg.
% mdparagraph:P105:end

\end{minipage}\par\addvspace{4pt}

\par\begin{minipage}{\linewidth}\raggedright
% mdparagraph:P106:begin
\textbf{F5.}
% mdparagraph:P106:end

% item: NEW-U05-F05
% mdparagraph:P107:begin
Give the selected indicative forms of cweðan: present 2sg and 3sg; preterite 3pl.
% mdparagraph:P107:end

\end{minipage}\par\addvspace{4pt}

\par\begin{minipage}{\linewidth}\raggedright
% mdparagraph:P108:begin
\textbf{F6.}
% mdparagraph:P108:end

% item: NEW-U05-F06
% mdparagraph:P109:begin
Give the selected indicative forms of giefan: present 2sg and 3sg; preterite 1sg and 1pl.
% mdparagraph:P109:end

\end{minipage}\par\addvspace{4pt}

\par\begin{minipage}{\linewidth}\raggedright
% mdparagraph:P110:begin
\textbf{F7.}
% mdparagraph:P110:end

% item: NEW-U05-F07
% mdparagraph:P111:begin
Give the selected indicative forms of biddan: present 1sg and 3sg; preterite 2sg.
% mdparagraph:P111:end

\end{minipage}\par\addvspace{4pt}

\par\begin{minipage}{\linewidth}\raggedright
% mdparagraph:P112:begin
\textbf{F8.}
% mdparagraph:P112:end

% item: NEW-U05-F08
% mdparagraph:P113:begin
Give the selected indicative forms of licgan: present 2sg and 3sg; preterite 1pl.
% mdparagraph:P113:end

\end{minipage}\par\addvspace{4pt}

\par\begin{minipage}{\linewidth}\raggedright
% mdparagraph:P114:begin
\textbf{F9.}
% mdparagraph:P114:end

% item: NEW-U05-F09
% mdparagraph:P115:begin
Give the selected indicative forms of sittan: present 3sg; preterite 1sg and 2sg.
% mdparagraph:P115:end

\end{minipage}\par\addvspace{4pt}

\par\begin{minipage}{\linewidth}\raggedright
% mdparagraph:P116:begin
\textbf{F10.}
% mdparagraph:P116:end

% item: NEW-U05-F10
% mdparagraph:P117:begin
Give the selected indicative forms of etan: present 2sg and 3sg; preterite 1sg and 2sg.
% mdparagraph:P117:end

\end{minipage}\par\addvspace{4pt}

\par\begin{minipage}{\linewidth}\raggedright
% mdparagraph:P118:begin
\textbf{F11.}
% mdparagraph:P118:end

% item: NEW-U05-F11
% mdparagraph:P119:begin
Give the selected indicative forms of sēon: present 1sg, 2sg and 3pl.
% mdparagraph:P119:end

\end{minipage}\par\addvspace{4pt}

\par\begin{minipage}{\linewidth}\raggedright
% mdparagraph:P120:begin
\textbf{F12.}
% mdparagraph:P120:end

% item: NEW-U05-F12
% mdparagraph:P121:begin
Give the selected indicative forms of gesēon: present 3sg; preterite 2sg and 2pl.
% mdparagraph:P121:end

\end{minipage}\par\addvspace{4pt}

\par\begin{minipage}{\linewidth}\raggedright
% mdparagraph:P122:begin
\textbf{F13.}
% mdparagraph:P122:end

% item: NEW-U05-F13
% mdparagraph:P123:begin
Give the selected indicative forms of faran: present 2sg and 3sg; preterite 1pl.
% mdparagraph:P123:end

\end{minipage}\par\addvspace{4pt}

\par\begin{minipage}{\linewidth}\raggedright
% mdparagraph:P124:begin
\textbf{F14.}
% mdparagraph:P124:end

% item: NEW-U05-F14
% mdparagraph:P125:begin
Give the selected indicative forms of slēan: present 1sg, 2sg and 2pl; preterite 2sg.
% mdparagraph:P125:end

\end{minipage}\par\addvspace{4pt}

\par\begin{minipage}{\linewidth}\raggedright
% mdparagraph:P126:begin
\textbf{F15.}
% mdparagraph:P126:end

% item: NEW-U05-F15
% mdparagraph:P127:begin
Give the selected indicative forms of standan: present 2sg and 3sg; preterite 3sg.
% mdparagraph:P127:end

\end{minipage}\par\addvspace{4pt}

\par\begin{minipage}{\linewidth}\raggedright
% mdparagraph:P128:begin
\textbf{F16.}
% mdparagraph:P128:end

% item: NEW-U05-F16
% mdparagraph:P129:begin
Give the selected indicative forms of hębban: present 1sg and 3sg; preterite 1pl.
% mdparagraph:P129:end

\end{minipage}\par\addvspace{4pt}

\par\begin{minipage}{\linewidth}\raggedright
% mdparagraph:P130:begin
\textbf{F17.}
% mdparagraph:P130:end

% item: NEW-U05-F17
% mdparagraph:P131:begin
Give the selected indicative forms of scieppan and gescieppan: scieppan: present 1sg and 2sg, preterite 3sg; gescieppan: preterite 3pl.
% mdparagraph:P131:end

\end{minipage}\par\addvspace{4pt}

\par\begin{minipage}{\linewidth}\raggedright
% mdparagraph:P132:begin
\textbf{F18.}
% mdparagraph:P132:end

% item: NEW-U05-F18
% mdparagraph:P133:begin
Give the selected indicative forms of swęrian: present 1sg, 2sg and 3sg; preterite 2sg.
% mdparagraph:P133:end

\end{minipage}\par\addvspace{4pt}

\par\begin{minipage}{\linewidth}\raggedright
% mdparagraph:P134:begin
\textbf{F19.}
% mdparagraph:P134:end

% item: NEW-U05-F19
% mdparagraph:P135:begin
Give the selected indicative forms of feallan, healdan and weaxan: present 3sg of each; preterite 1pl of weaxan.
% mdparagraph:P135:end

\end{minipage}\par\addvspace{4pt}

\par\begin{minipage}{\linewidth}\raggedright
% mdparagraph:P136:begin
\textbf{F20.}
% mdparagraph:P136:end

% item: NEW-U05-F20
% mdparagraph:P137:begin
Give the selected indicative forms of cnāwan and wēpan: present 2sg and preterite 2sg of each.
% mdparagraph:P137:end

\end{minipage}\par\addvspace{4pt}

\par\begin{minipage}{\linewidth}\raggedright
% mdparagraph:P138:begin
\textbf{F21.}
% mdparagraph:P138:end

% item: NEW-U05-F21
% mdparagraph:P139:begin
Give the selected indicative forms of hēawan and flōwan: present 2sg and preterite 2sg of each.
% mdparagraph:P139:end

\end{minipage}\par\addvspace{4pt}

\par\begin{minipage}{\linewidth}\raggedright
% mdparagraph:P140:begin
\textbf{F22.}
% mdparagraph:P140:end

% item: NEW-U05-F22
% mdparagraph:P141:begin
Give the selected indicative forms of lǣtan, hātan and fōn: present 3sg of each; preterite 2sg of fōn.
% mdparagraph:P141:end

\end{minipage}\par\addvspace{4pt}

\par\begin{minipage}{\linewidth}\raggedright
% mdparagraph:P142:begin
\textbf{F23.}
% mdparagraph:P142:end

% item: NEW-U05-F23
% mdparagraph:P143:begin
Give the selected indicative forms of geswīcan and scīnan: present 3sg and preterite 2sg of each.
% mdparagraph:P143:end

\end{minipage}\par\addvspace{4pt}

\par\begin{minipage}{\linewidth}\raggedright
% mdparagraph:P144:begin
\textbf{F24.}
% mdparagraph:P144:end

% item: NEW-U05-F24
% mdparagraph:P145:begin
Give the selected indicative forms of flēogan and swimman: present 3sg and preterite 3pl of each.
% mdparagraph:P145:end

\end{minipage}\par\addvspace{4pt}

\reserve{7\baselineskip}
% mdparagraph:P146:begin
\subsection*{Clauses — C1–C8}
% mdparagraph:P146:end

\reserve{5\baselineskip}
\par\begin{minipage}{\linewidth}\raggedright
% mdparagraph:P147:begin
\textbf{C1.}
% mdparagraph:P147:end

% item: NEW-U05-C01
% mdparagraph:P148:begin
Translate: You (singular) carry the food, and I take the fish.
% mdparagraph:P148:end

\end{minipage}\par\addvspace{4pt}

\par\begin{minipage}{\linewidth}\raggedright
% mdparagraph:P149:begin
\textbf{C2.}
% mdparagraph:P149:end

% item: NEW-U05-C02
% mdparagraph:P150:begin
Translate using giefan: We (two) gave the men our food (belonging to us two).
% mdparagraph:P150:end

\end{minipage}\par\addvspace{4pt}

\par\begin{minipage}{\linewidth}\raggedright
% mdparagraph:P151:begin
\textbf{C3.}
% mdparagraph:P151:end

% item: NEW-U05-C03
% mdparagraph:P152:begin
Translate: The companions sat in the dwelling, and the man stood on the shore.
% mdparagraph:P152:end

\end{minipage}\par\addvspace{4pt}

\par\begin{minipage}{\linewidth}\raggedright
% mdparagraph:P153:begin
\textbf{C4.}
% mdparagraph:P153:end

% item: NEW-U05-C04
% mdparagraph:P154:begin
Translate using etan: I did not eat the food; you (singular) ate it.
% mdparagraph:P154:end

\end{minipage}\par\addvspace{4pt}

\par\begin{minipage}{\linewidth}\raggedright
% mdparagraph:P155:begin
\textbf{C5.}
% mdparagraph:P155:end

% item: NEW-U05-C05
% mdparagraph:P156:begin
Rewrite with hīe as subject and plural objects: Hē geseah ðone fugol and nōm ðone fisc.
% mdparagraph:P156:end

\end{minipage}\par\addvspace{4pt}

\par\begin{minipage}{\linewidth}\raggedright
% mdparagraph:P157:begin
\textbf{C6.}
% mdparagraph:P157:end

% item: NEW-U05-C06
% mdparagraph:P158:begin
Translate using cweðan to introduce direct speech: The man says, “I see the trees.”
% mdparagraph:P158:end

\end{minipage}\par\addvspace{4pt}

\par\begin{minipage}{\linewidth}\raggedright
% mdparagraph:P159:begin
\textbf{C7.}
% mdparagraph:P159:end

% item: NEW-U05-C07
% mdparagraph:P160:begin
Translate using geswīcan: You (plural) cease from your work.
% mdparagraph:P160:end

\end{minipage}\par\addvspace{4pt}

\par\begin{minipage}{\linewidth}\raggedright
% mdparagraph:P161:begin
\textbf{C8.}
% mdparagraph:P161:end

% item: NEW-U05-C08
% mdparagraph:P162:begin
Translate: The tree falls, and the men hold the ship.
% mdparagraph:P162:end

\end{minipage}\par\addvspace{4pt}

\reserve{7\baselineskip}
% mdparagraph:P163:begin
\section*{9. Cook I.5}
% mdparagraph:P163:end

% mdparagraph:P164:begin
Translate all sixteen sentences into Old English. These are Cook’s unchanged English prompts from \textbf{Exercises in Old English}, Group I, Exercise 5, printed p.3 (one-based PDF p.11). Use the ordinary constructions taught above and the earlier weak-verb and nominal systems. Original item numbers are retained.
% mdparagraph:P164:end

\reserve{5\baselineskip}
\par\begin{minipage}{\linewidth}\raggedright
% mdparagraph:P165:begin
\textbf{I.5.1.}
% mdparagraph:P165:end

% item: EX-I-05-01
% mdparagraph:P166:begin
God created the heavens and the earth.
% mdparagraph:P166:end

\end{minipage}\par\addvspace{4pt}

\par\begin{minipage}{\linewidth}\raggedright
% mdparagraph:P167:begin
\textbf{I.5.2.}
% mdparagraph:P167:end

% item: EX-I-05-02
% mdparagraph:P168:begin
God made two great lights.
% mdparagraph:P168:end

\end{minipage}\par\addvspace{4pt}

\par\begin{minipage}{\linewidth}\raggedright
% mdparagraph:P169:begin
\textbf{I.5.3.}
% mdparagraph:P169:end

% item: EX-I-05-03
% mdparagraph:P170:begin
God made the firmament.
% mdparagraph:P170:end

\end{minipage}\par\addvspace{4pt}

\par\begin{minipage}{\linewidth}\raggedright
% mdparagraph:P171:begin
\textbf{I.5.4.}
% mdparagraph:P171:end

% item: EX-I-05-04
% mdparagraph:P172:begin
God created the great whales.
% mdparagraph:P172:end

\end{minipage}\par\addvspace{4pt}

\par\begin{minipage}{\linewidth}\raggedright
% mdparagraph:P173:begin
\textbf{I.5.5.}
% mdparagraph:P173:end

% item: EX-I-05-05
% mdparagraph:P174:begin
He made the beasts of the earth.
% mdparagraph:P174:end

\end{minipage}\par\addvspace{4pt}

% teaching: NEW-U05-GRAM-U05-SOURCE-ALL-THOSE-THINGS
% teaching: NEW-U05-SUPFORM-eall-acc-n-pl-strong
% teaching: NEW-U05-SUPSENSE-eall-all
% mdparagraph:P175:begin
\textbf{Supplied phrase for I.5.6:} eall ðā ðing — all those things (accusative plural); use the supplied phrase. Here eall is strong neuter accusative plural, ðā is demonstrative accusative plural, and ðing is neuter accusative plural. The whole phrase remains available; neither eall nor the quantifier construction requires independent production.
% mdparagraph:P175:end

\reserve{5\baselineskip}
\par\begin{minipage}{\linewidth}\raggedright
% mdparagraph:P176:begin
\textbf{I.5.6.}
% mdparagraph:P176:end

% item: EX-I-05-06
% mdparagraph:P177:begin
He saw all those things.
% mdparagraph:P177:end

\end{minipage}\par\addvspace{4pt}

\par\begin{minipage}{\linewidth}\raggedright
% mdparagraph:P178:begin
\textbf{I.5.7.}
% mdparagraph:P178:end

% item: EX-I-05-07
% mdparagraph:P179:begin
He finished his work.
% mdparagraph:P179:end

\end{minipage}\par\addvspace{4pt}

\par\begin{minipage}{\linewidth}\raggedright
% mdparagraph:P180:begin
\textbf{I.5.8.}
% mdparagraph:P180:end

% item: EX-I-05-08
% mdparagraph:P181:begin
He blesses the seventh day.
% mdparagraph:P181:end

\end{minipage}\par\addvspace{4pt}

\par\begin{minipage}{\linewidth}\raggedright
% mdparagraph:P182:begin
\textbf{I.5.9.}
% mdparagraph:P182:end

% item: EX-I-05-09
% mdparagraph:P183:begin
The man ceases from his work.
% mdparagraph:P183:end

\end{minipage}\par\addvspace{4pt}

\par\begin{minipage}{\linewidth}\raggedright
% mdparagraph:P184:begin
\textbf{I.5.10.}
% mdparagraph:P184:end

% item: EX-I-05-10
% mdparagraph:P185:begin
I see the bird in the tree.
% mdparagraph:P185:end

\end{minipage}\par\addvspace{4pt}

\par\begin{minipage}{\linewidth}\raggedright
% mdparagraph:P186:begin
\textbf{I.5.11.}
% mdparagraph:P186:end

% item: EX-I-05-11
% mdparagraph:P187:begin
They saw the fish in the water.
% mdparagraph:P187:end

\end{minipage}\par\addvspace{4pt}

\par\begin{minipage}{\linewidth}\raggedright
% mdparagraph:P188:begin
\textbf{I.5.12.}
% mdparagraph:P188:end

% item: EX-I-05-12
% mdparagraph:P189:begin
The birds of the air fly.
% mdparagraph:P189:end

\end{minipage}\par\addvspace{4pt}

\par\begin{minipage}{\linewidth}\raggedright
% mdparagraph:P190:begin
\textbf{I.5.13.}
% mdparagraph:P190:end

% item: EX-I-05-13
% mdparagraph:P191:begin
The trees grew.
% mdparagraph:P191:end

\end{minipage}\par\addvspace{4pt}

\par\begin{minipage}{\linewidth}\raggedright
% mdparagraph:P192:begin
\textbf{I.5.14.}
% mdparagraph:P192:end

% item: EX-I-05-14
% mdparagraph:P193:begin
The stars shone.
% mdparagraph:P193:end

\end{minipage}\par\addvspace{4pt}

\par\begin{minipage}{\linewidth}\raggedright
% mdparagraph:P194:begin
\textbf{I.5.15.}
% mdparagraph:P194:end

% item: EX-I-05-15
% mdparagraph:P195:begin
The fish swims.
% mdparagraph:P195:end

\end{minipage}\par\addvspace{4pt}

\par\begin{minipage}{\linewidth}\raggedright
% mdparagraph:P196:begin
\textbf{I.5.16.}
% mdparagraph:P196:end

% item: EX-I-05-16
% mdparagraph:P197:begin
The men swam.
% mdparagraph:P197:end

\end{minipage}\par\addvspace{4pt}

\reserve{7\baselineskip}
% mdparagraph:P198:begin
\section*{10. Connected reading: Food for the companions and the birds}
% mdparagraph:P198:end

% mdparagraph:P199:begin
This account is new teaching material. The preterites narrate successive actions. Read the pronouns in context: hīe after the giving clause can naturally refer to the companions, but its form alone does not force that antecedent. Ān mann does not say whether the man belongs to the original pair or their companions.
% mdparagraph:P199:end

% composition: NEW-U05-TEXT-A
% mdparagraph:P200:begin
\par\begin{minipage}{\linewidth}\raggedright
\readingtext{Twēgen męnn bǣron męte. Hiera gefēran sǣton on ðǣm hofe. Ðā męnn gēafon ðǣm gefērum ðone męte, and hīe ǣton. Ān mann ne ǣt. Hē geaf his męte ðǣm fuglum. Ðā fuglas ǣton and flugon. Se mann sæt on ðǣm strande and geseah ðā fiscas. Ðā fiscas swummon on ðǣm wætere.}
\end{minipage}\par
% mdparagraph:P200:end

\reserve{7\baselineskip}
% mdparagraph:P201:begin
\subsection*{Connected reading — T1–T4}
% mdparagraph:P201:end

\reserve{5\baselineskip}
\par\begin{minipage}{\linewidth}\raggedright
% mdparagraph:P202:begin
\textbf{T1.}
% mdparagraph:P202:end

% item: NEW-U05-T01
% mdparagraph:P203:begin
In the new account, what did one man do with his food, and what did its recipients do next?
% mdparagraph:P203:end

\end{minipage}\par\addvspace{4pt}

\par\begin{minipage}{\linewidth}\raggedright
% mdparagraph:P204:begin
\textbf{T2.}
% mdparagraph:P204:end

% item: NEW-U05-T02
% mdparagraph:P205:begin
In ðǣm gefērum ðone męte, identify the recipient and the thing given, and explain their cases. What subject controls gēafon?
% mdparagraph:P205:end

\end{minipage}\par\addvspace{4pt}

\par\begin{minipage}{\linewidth}\raggedright
% mdparagraph:P206:begin
\textbf{T3.}
% mdparagraph:P206:end

% item: NEW-U05-T03
% mdparagraph:P207:begin
Retell the last two sentences of the new account in the present indicative. Keep the subjects, objects and location phrases.
% mdparagraph:P207:end

\end{minipage}\par\addvspace{4pt}

\par\begin{minipage}{\linewidth}\raggedright
% mdparagraph:P208:begin
\textbf{T4.}
% mdparagraph:P208:end

% item: NEW-U05-T04
% mdparagraph:P209:begin
Translate the complete new account.
% mdparagraph:P209:end

\end{minipage}\par\addvspace{4pt}

\reserve{7\baselineskip}
% mdparagraph:P210:begin
\section*{11. Source reading: The first day revisited}
% mdparagraph:P210:end

% mdparagraph:P211:begin
This is the same complete first-day passage read in Unit 1, now revisited with the strong finite forms learned in this unit. The original text and assignment are unchanged. Use all the local help below for the constructions that remain supplied. Cook prints the passage in \textbf{A First Book in Old English}, Reader I, pp.123–124 (one-based PDF pp.139–140).
% mdparagraph:P211:end

% source: FB-I-CREATION-FIRST-DAY-01
% mdparagraph:P212:begin
\par\begin{minipage}{\linewidth}\raggedright
\readingtext{On anginne gescēop God heofonan and eorðan. Sēo eorðe sōðlice wæs īdel and ǣmtig; and ðīestru wǣron ofer ðǣre nēowolnesse brādnesse; and Godes gāst wæs gefęred ofer wæteru. God cwæð ðā, “Geweorðe lēoht”; and lēoht wearð geworht. God geseah ðā ðæt hit gōd wæs; and hē gedǣlde ðæt lēoht fram ðǣm ðīestrum. And hēt ðæt lēoht Dæg, and ðā ðīestru Niht. Ðā wæs geworden ǣfen and morgen ān dæg.}
\end{minipage}\par
% mdparagraph:P212:end

\reserve{7\baselineskip}
% mdparagraph:P213:begin
\subsection*{Local help for the first day}
% mdparagraph:P213:end

% mdparagraph:P214:begin
Each explanation applies to the actual printed phrase and remains available during R1–R4. Learned indicative morphology does not make these other constructions independent production targets.
% mdparagraph:P214:end

\reserve{5\baselineskip}
% mdparagraph:P215:begin
\textbf{On anginne}
% mdparagraph:P215:end

% teaching: NEW-U05-SUPSENSE-on-time-supplied
% teaching: NEW-U05-GRAM-READING-ON-ANGINNE
% teaching: NEW-U05-SUPPORT-U05-SRC-D1-ON-BEGINNING
% mdparagraph:P216:begin
On anginne: in the beginning.
% mdparagraph:P216:end

% mdparagraph:P217:begin
This is temporal on with a dative. The whole time phrase is supplied; the independent on-construction already practised states location.
% mdparagraph:P217:end

\reserve{5\baselineskip}
% mdparagraph:P218:begin
\textbf{heofonan}
% mdparagraph:P218:end

% teaching: NEW-U05-SUPFORM-heofone-acc-sg
% teaching: NEW-U05-SUPSENSE-heofone-1
% teaching: NEW-U05-GRAM-READING-HEOFONE
% teaching: NEW-U05-SUPPORT-U05-SRC-D1-HEAVEN
% mdparagraph:P219:begin
heofonan: heaven/\allowbreak{}the sky, accusative singular of weak feminine heofone.
% mdparagraph:P219:end

% mdparagraph:P220:begin
Heofone is a weak feminine noun, and this exact heofonan is accusative singular. It is distinct from the learned masculine heofon and its plural heofonas. No further forms of this feminine noun are required here.
% mdparagraph:P220:end

\reserve{5\baselineskip}
% mdparagraph:P221:begin
\textbf{sōðlice}
% mdparagraph:P221:end

% teaching: NEW-U05-SUPSENSE-sothlice-1
% teaching: NEW-U05-GRAM-U05-SOURCE-CONNECTIVE-SOTHLICE
% teaching: NEW-U05-SUPPORT-U05-SRC-D1-CONNECTIVE
% mdparagraph:P222:begin
sōðlice: now; for its part.
% mdparagraph:P222:end

% mdparagraph:P223:begin
This connective use differs from the earlier ordinary meaning “truly.” Cook’s note gives Latin autem. Keep the contextual meaning supplied here.
% mdparagraph:P223:end

\reserve{5\baselineskip}
% mdparagraph:P224:begin
\textbf{ofer ðǣre nēowolnesse brādnesse;}
% mdparagraph:P224:end

% teaching: NEW-U05-SUPFORM-ofer-invariant
% teaching: NEW-U05-SUPSENSE-ofer-1
% teaching: NEW-U05-GRAM-READING-OFER-DATIVE
% teaching: NEW-U05-SUPPORT-U05-SRC-D1-OFER-DEEP
% mdparagraph:P225:begin
ofer ðǣre nēowolnesse brādnesse: over the surface of the deep.
% mdparagraph:P225:end

% mdparagraph:P226:begin
Ofer governs dative brādnesse. The preceding ðǣre nēowolnesse is feminine genitive singular, dependent on brādnesse. Learn the analysis of this exact phrase without extending it to a complete rule for ofer.
% mdparagraph:P226:end

\reserve{5\baselineskip}
% mdparagraph:P227:begin
\textbf{wæs gefęred ofer wæteru.}
% mdparagraph:P227:end

% teaching: NEW-U05-SUPFORM-geferian-past-participle
% teaching: NEW-U05-SUPSENSE-geferian-1
% teaching: NEW-U05-GRAM-U05-SOURCE-CARRIED-OVER-WATERS
% teaching: NEW-U05-SUPPORT-U05-SRC-D1-CARRIED
% mdparagraph:P228:begin
wæs gefęred ofer wæteru: was borne/\allowbreak{}carried (moved) over the waters.
% mdparagraph:P228:end

% mdparagraph:P229:begin
Wæs is finite; gefęred is the supplied participle in a passive expression. Here ofer has accusative plural wæteru. Both the passive and this local prepositional construction remain supplied.
% mdparagraph:P229:end

\reserve{5\baselineskip}
% mdparagraph:P230:begin
\textbf{ðā /\allowbreak{} Ðā}
% mdparagraph:P230:end

% teaching: NEW-U05-SUPFORM-tha-adverb-invariant
% teaching: NEW-U05-SUPSENSE-tha-adverb-1
% teaching: NEW-U05-GRAM-READING-TEMPORAL-THA
% teaching: NEW-U05-SUPPORT-U05-SRC-D1-THA
% mdparagraph:P231:begin
ðā /\allowbreak{} Ðā: then.
% mdparagraph:P231:end

% mdparagraph:P232:begin
These occurrences are temporal adverbs. They are distinct from the identically spelled plural article or demonstrative.
% mdparagraph:P232:end

\reserve{5\baselineskip}
% mdparagraph:P233:begin
\textbf{“Geweorðe lēoht”;}
% mdparagraph:P233:end

% teaching: NEW-U05-SUPFORM-geweorthan-pres-subj-3-sg
% teaching: NEW-U05-SUPSENSE-geweorthan-1
% teaching: NEW-U05-GRAM-READING-OPTATIVE-COMMAND
% teaching: NEW-U05-SUPPORT-U05-SRC-D1-WISH
% mdparagraph:P234:begin
Geweorðe lēoht: Let light come to be.
% mdparagraph:P234:end

% mdparagraph:P235:begin
Geweorðe is a third-singular present subjunctive used as an optative, expressing a wish or command. The whole expression is supplied; it is outside the indicative tables.
% mdparagraph:P235:end

\reserve{5\baselineskip}
% mdparagraph:P236:begin
\textbf{wearð geworht.}
% mdparagraph:P236:end

% teaching: NEW-U05-SUPFORM-gewyrcean-past-participle
% teaching: NEW-U05-SUPSENSE-weorthan-1
% teaching: NEW-U05-GRAM-READING-PASSIVE-GEWORHT
% teaching: NEW-U05-SUPPORT-U05-SRC-D1-MADE
% mdparagraph:P237:begin
wearð geworht: was made; came into being.
% mdparagraph:P237:end

% mdparagraph:P238:begin
Wearð is the finite preterite. The ordinary meaning of weorðan, “become,” was learned earlier; here wearð serves with the supplied participle geworht in a passive expression. That auxiliary use remains supplied.
% mdparagraph:P238:end

\reserve{5\baselineskip}
% mdparagraph:P239:begin
\textbf{geseah ðā ðæt hit gōd wæs;}
% mdparagraph:P239:end

% teaching: NEW-U05-SUPFORM-thaet-conjunction-invariant
% teaching: NEW-U05-SUPSENSE-thaet-conjunction-1
% teaching: NEW-U05-GRAM-READING-CONTENT-CLAUSE
% teaching: NEW-U05-SUPPORT-U05-SRC-D1-CONTENT
% mdparagraph:P240:begin
God geseah … ðæt hit gōd wæs: God saw … that it was good.
% mdparagraph:P240:end

% mdparagraph:P241:begin
Geseah is now a learned indicative form. The conjunction ðæt introduces the supplied content clause; its final wæs belongs inside that clause. Hit refers to neuter lēoht. This ðæt means “that,” not “the.”
% mdparagraph:P241:end

\reserve{5\baselineskip}
% mdparagraph:P242:begin
\textbf{gedǣlde ðæt lēoht fram ðǣm ðīestrum.}
% mdparagraph:P242:end

% teaching: NEW-U05-SUPFORM-gedaelan-pret-3-sg
% teaching: NEW-U05-SUPSENSE-gedaelan-1
% teaching: NEW-U05-GRAM-U05-SOURCE-GEDAELDE
% teaching: NEW-U05-SUPPORT-U05-SRC-D1-DIVIDE
% mdparagraph:P243:begin
gedǣlde ðæt lēoht fram ðǣm ðīestrum: separated the light from the darkness.
% mdparagraph:P243:end

% mdparagraph:P244:begin
Gedǣlde is the supplied weak third-singular preterite of gedǣlan, “divide; separate.” The separating fram + dative construction is already familiar; this additional verb form remains supplied.
% mdparagraph:P244:end

\reserve{5\baselineskip}
% mdparagraph:P245:begin
\textbf{hēt ðæt lēoht Dæg, and ðā ðīestru Niht.}
% mdparagraph:P245:end

% teaching: NEW-U05-GRAM-READING-NAMING
% teaching: NEW-U05-SUPPORT-U05-SRC-D1-NAMING
% mdparagraph:P246:begin
He called the light Day and the darkness Night.
% mdparagraph:P246:end

% mdparagraph:P247:begin
The object and the name given to it are both accusative: lēoht/\allowbreak{}Dæg and ðīestru/\allowbreak{}Niht. The continuing subject is God; the second naming clause also leaves the verb understood. The finite hēt is learned, while this two-accusative naming construction remains supplied.
% mdparagraph:P247:end

\reserve{5\baselineskip}
% mdparagraph:P248:begin
\textbf{Ðā wæs geworden ǣfen and morgen ān dæg.}
% mdparagraph:P248:end

% teaching: NEW-U05-SUPFORM-geweorthan-past-participle
% teaching: NEW-U05-GRAM-READING-RESULTATIVE-ENDING
% teaching: NEW-U05-SUPPORT-U05-SRC-D1-CLOSE
% mdparagraph:P249:begin
Then evening and morning had come to be—one day.
% mdparagraph:P249:end

% mdparagraph:P250:begin
Wæs geworden is a supplied resultative/\allowbreak{}perfect expression. It describes what has come to be, rather than the passive making in wearð geworht. Preserve ān dæg as “one day.” The formula is given whole; no unique subject/\allowbreak{}predicate division or new agreement rule is required.
% mdparagraph:P250:end

\reserve{7\baselineskip}
% mdparagraph:P251:begin
\subsection*{First-day revisit — R1–R4}
% mdparagraph:P251:end

\reserve{5\baselineskip}
\par\begin{minipage}{\linewidth}\raggedright
% mdparagraph:P252:begin
\textbf{R1.}
% mdparagraph:P252:end

% item: NEW-U05-R01
% mdparagraph:P253:begin
In the first-day text, identify the lemma, strong class, tense, person and number of gescēop, cwæð, geseah and hēt. Use the clauses to identify the subject of hēt.
% mdparagraph:P253:end

\end{minipage}\par\addvspace{4pt}

% teaching: NEW-U05-CIT-weorthan
% mdparagraph:P254:begin
\textbf{Supplied citation for R2:} weorðan, “become.” The headword is printed for identification; no infinitival construction or unaided citation recall is required.
% mdparagraph:P254:end

\reserve{5\baselineskip}
\par\begin{minipage}{\linewidth}\raggedright
% mdparagraph:P255:begin
\textbf{R2.}
% mdparagraph:P255:end

% item: NEW-U05-R02
% mdparagraph:P256:begin
Use the first-day notes to distinguish the ordinary meaning of weorðan from its use in wearð geworht. Which word is finite in that supplied phrase?
% mdparagraph:P256:end

\end{minipage}\par\addvspace{4pt}

\par\begin{minipage}{\linewidth}\raggedright
% mdparagraph:P257:begin
\textbf{R3.}
% mdparagraph:P257:end

% item: NEW-U05-R03
% mdparagraph:P258:begin
The first-day text has heofonan; Cook I.5.1 says the heavens. Using the source note, explain the difference in noun, gender and number.
% mdparagraph:P258:end

\end{minipage}\par\addvspace{4pt}

\par\begin{minipage}{\linewidth}\raggedright
% mdparagraph:P259:begin
\textbf{R4.}
% mdparagraph:P259:end

% item: NEW-U05-R04
% mdparagraph:P260:begin
Translate the complete first-day passage, using the printed local help where needed.
% mdparagraph:P260:end

\end{minipage}\par\addvspace{4pt}

\reserve{7\baselineskip}
% mdparagraph:P261:begin
\section*{12. Source reading: The second day}
% mdparagraph:P261:end

% mdparagraph:P262:begin
Continue with the complete second-day paragraph, Cook’s Reader I, p.124, editorial prose lines 4–10 (one-based PDF p.140). The edition prints marginal line anchors at 5 and 10; the intervening line count is editorial. The new finite knowledge helps with the narrative, while the wishes, relative clauses and closing expressions retain local help.
% mdparagraph:P262:end

% source: FB-I-CREATION-SECOND-DAY-01
% mdparagraph:P263:begin
\par\begin{minipage}{\linewidth}\raggedright
\readingtext{God cwæð ðā ęft, “Geweorðe nū fæstnes tōmiddes ðǣm wæterum, and tōtwǣme ðā wæteru fram ðǣm wæterum.” And God geworhte ðā fæstnesse, and tōtwǣmde ðā wæteru ðe wǣron under ðǣre fæstnesse fram ðǣm ðe wǣron bufan ðǣre fæstnesse; hit wæs ðā swā gedōn. And God hēt ðā fæstnesse Heofonan. And wæs ðā geworden ǣfen and morgen ōðer dæg.}
\end{minipage}\par
% mdparagraph:P263:end

\reserve{7\baselineskip}
% mdparagraph:P264:begin
\subsection*{Local help for the second day}
% mdparagraph:P264:end

% mdparagraph:P265:begin
The following six groups of help remain available during R5–R8. They explain the actual forms and their relations without requiring independent production of the supplied constructions.
% mdparagraph:P265:end

\reserve{5\baselineskip}
% mdparagraph:P266:begin
\textbf{ðā, ęft, nū, swā}
% mdparagraph:P266:end

% teaching: NEW-U05-SUPFORM-eft-invariant
% teaching: NEW-U05-SUPSENSE-eft-again-also
% teaching: NEW-U05-SUPFORM-nu-invariant
% teaching: NEW-U05-SUPSENSE-nu-1
% teaching: NEW-U05-SUPFORM-swa-invariant
% teaching: NEW-U05-SUPSENSE-swa-1
% teaching: NEW-U05-GRAM-U05-SOURCE-SECOND-DAY-ADVERBS
% teaching: NEW-U05-SUPPORT-U05-SRC-D2-ADVERBS
% mdparagraph:P267:begin
ðā: then; ęft: again/\allowbreak{}also; nū: now; swā: so/\allowbreak{}thus.
% mdparagraph:P267:end

% mdparagraph:P268:begin
These exact adverbial uses are supplied. Cook glosses ęft with Latin quoque, “also.” Preserve the printed ę.
% mdparagraph:P268:end

\reserve{5\baselineskip}
% mdparagraph:P269:begin
\textbf{“Geweorðe nū fæstnes tōmiddes ðǣm wæterum, and tōtwǣme ðā wæteru fram ðǣm wæterum.”}
% mdparagraph:P269:end

% teaching: NEW-U05-SUPFORM-tomiddes-invariant
% teaching: NEW-U05-SUPSENSE-tomiddes-midst
% teaching: NEW-U05-SUPFORM-totwaeman-pres-subj-3-sg
% teaching: NEW-U05-SUPSENSE-totwaeman-separate
% teaching: NEW-U05-GRAM-U05-SOURCE-SECOND-DAY-WISH
% teaching: NEW-U05-SUPPORT-U05-SRC-D2-WISH
% mdparagraph:P270:begin
Let a firmament come to be in the midst of the waters, and let it separate the waters from the waters.
% mdparagraph:P270:end

% mdparagraph:P271:begin
Geweorðe and tōtwǣme are coordinated third-singular present optatives; fæstnes is the subject of both. Tōmiddes has dative ðǣm wæterum. Tōtwǣme, “let it separate,” has accusative ðā wæteru and fram + dative. The whole wish and its new preposition remain supplied.
% mdparagraph:P271:end

\reserve{5\baselineskip}
% mdparagraph:P272:begin
\textbf{tōtwǣmde ðā wæteru ðe wǣron under ðǣre fæstnesse fram ðǣm ðe wǣron bufan ðǣre fæstnesse;}
% mdparagraph:P272:end

% teaching: NEW-U05-SUPFORM-totwaeman-pret-3-sg
% teaching: NEW-U05-SUPFORM-the-relative-subject-invariant
% teaching: NEW-U05-SUPSENSE-the-relative-subject-waters
% teaching: NEW-U05-SUPFORM-under-invariant
% teaching: NEW-U05-SUPSENSE-under-below
% teaching: NEW-U05-SUPFORM-bufan-invariant
% teaching: NEW-U05-SUPSENSE-bufan-above
% teaching: NEW-U05-SUPSENSE-se-those-waters-substantival
% teaching: NEW-U05-GRAM-U05-SOURCE-SECOND-DAY-RELATIVES
% teaching: NEW-U05-SUPPORT-U05-SRC-D2-SEPARATING-RELATIVES
% mdparagraph:P273:begin
and separated the waters which were under the firmament from those which were above the firmament.
% mdparagraph:P273:end

% mdparagraph:P274:begin
Tōtwǣmde is the supplied third-singular weak preterite, “separated.” In each relative clause, ðe is the subject of wǣron. The first refers to the waters below the firmament. After fram, dative ðǣm means “those (waters)”; its following relative describes the waters above. Under and bufan both govern dative ðǣre fæstnesse here. These subject-relative clauses and the omitted noun after “those” remain supplied.
% mdparagraph:P274:end

\reserve{5\baselineskip}
% mdparagraph:P275:begin
\textbf{hit wæs ðā swā gedōn.}
% mdparagraph:P275:end

% teaching: NEW-U05-SUPFORM-gedon-past-participle
% teaching: NEW-U05-SUPSENSE-gedon-do-make
% teaching: NEW-U05-GRAM-U05-SOURCE-SECOND-DAY-FULFILLMENT
% teaching: NEW-U05-SUPPORT-U05-SRC-D2-FULFILLMENT
% mdparagraph:P276:begin
hit wæs ðā swā gedōn: it was then done/\allowbreak{}made so.
% mdparagraph:P276:end

% mdparagraph:P277:begin
Wæs is finite; gedōn is the supplied past participle of gedōn, “do; make.” The whole passive/\allowbreak{}resultative expression describes the fulfillment of the preceding action.
% mdparagraph:P277:end

\reserve{5\baselineskip}
% mdparagraph:P278:begin
\textbf{God hēt ðā fæstnesse Heofonan.}
% mdparagraph:P278:end

% teaching: NEW-U05-GRAM-U05-SOURCE-SECOND-DAY-NAMING
% teaching: NEW-U05-SUPPORT-U05-SRC-D2-NAMING
% mdparagraph:P279:begin
God called the firmament Heaven.
% mdparagraph:P279:end

% mdparagraph:P280:begin
Fæstnesse is the accusative singular object; Heofonan is the accusative singular name, from weak feminine heofone. The hēt form is learned, while this naming construction and feminine noun form remain supplied.
% mdparagraph:P280:end

\reserve{5\baselineskip}
% mdparagraph:P281:begin
\textbf{And wæs ðā geworden ǣfen and morgen ōðer dæg.}
% mdparagraph:P281:end

% teaching: NEW-U05-GRAM-U05-SOURCE-SECOND-DAY-CLOSE
% teaching: NEW-U05-SUPPORT-U05-SRC-D2-CLOSE
% mdparagraph:P282:begin
And evening and morning had come to be—a second day.
% mdparagraph:P282:end

% mdparagraph:P283:begin
The resultative closing formula remains supplied. Here ōðer means “second” and is masculine nominative singular, agreeing with dæg. Contrast the first-day formula’s literal “one day.”
% mdparagraph:P283:end

\reserve{7\baselineskip}
% mdparagraph:P284:begin
\subsection*{Second day — R5–R8}
% mdparagraph:P284:end

\reserve{5\baselineskip}
\par\begin{minipage}{\linewidth}\raggedright
% mdparagraph:P285:begin
\textbf{R5.}
% mdparagraph:P285:end

% item: NEW-U05-R05
% mdparagraph:P286:begin
In the second-day text, distinguish cwæð and hēt from geworhte and the supplied tōtwǣmde. Identify strong and weak preterites, and identify their common subject.
% mdparagraph:P286:end

\end{minipage}\par\addvspace{4pt}

\par\begin{minipage}{\linewidth}\raggedright
% mdparagraph:P287:begin
\textbf{R6.}
% mdparagraph:P287:end

% item: NEW-U05-R06
% mdparagraph:P288:begin
Using the second-day local help, explain which waters are separated from which. What does ðǣm refer to in fram ðǣm ðe wǣron bufan ðǣre fæstnesse?
% mdparagraph:P288:end

\end{minipage}\par\addvspace{4pt}

\par\begin{minipage}{\linewidth}\raggedright
% mdparagraph:P289:begin
\textbf{R7.}
% mdparagraph:P289:end

% item: NEW-U05-R07
% mdparagraph:P290:begin
What does God call the firmament in the second-day account? Explain why ōðer agrees with dæg in the closing formula, while keeping the formula itself supplied.
% mdparagraph:P290:end

\end{minipage}\par\addvspace{4pt}

\par\begin{minipage}{\linewidth}\raggedright
% mdparagraph:P291:begin
\textbf{R8.}
% mdparagraph:P291:end

% item: NEW-U05-R08
% mdparagraph:P292:begin
Translate the complete second-day passage, using the printed local help where needed.
% mdparagraph:P292:end

\end{minipage}\par\addvspace{4pt}

\clearpage
% mdparagraph:P293:begin
\section*{13. Cumulative review}
% mdparagraph:P293:end

% mdparagraph:P294:begin
These three exercises are new teaching material. The printed citations in X2 remain supplied; compare the finite systems without adding nonfinite constructions.
% mdparagraph:P294:end

\reserve{5\baselineskip}
\par\begin{minipage}{\linewidth}\raggedright
% mdparagraph:P295:begin
\textbf{X1.}
% mdparagraph:P295:end

% item: NEW-U05-X01
% mdparagraph:P296:begin
Compare hē geaf and hē sealde. Give the corresponding preterite clauses with ðū, and explain why the endings differ.
% mdparagraph:P296:end

\end{minipage}\par\addvspace{4pt}

% teaching: NEW-U05-CIT-lufian
% mdparagraph:P297:begin
\textbf{Supplied citations for X2:} lufian and (ge)lufod. The parentheses in (ge)lufod preserve the earlier source citation notation; they do not introduce a rule for prefixes, perfect or passive syntax.
% mdparagraph:P297:end

\reserve{5\baselineskip}
\par\begin{minipage}{\linewidth}\raggedright
% mdparagraph:P298:begin
\textbf{X2.}
% mdparagraph:P298:end

% item: NEW-U05-X02
% mdparagraph:P299:begin
The printed entries give swęrian, swōr, swōron, sworen and lufian, lufode, (ge)lufod. Explain why -ian alone does not identify a weak verb. Give the present first-person singular of each.
% mdparagraph:P299:end

\end{minipage}\par\addvspace{4pt}

\par\begin{minipage}{\linewidth}\raggedright
% mdparagraph:P300:begin
\textbf{X3.}
% mdparagraph:P300:end

% item: NEW-U05-X03
% mdparagraph:P301:begin
Within the taught indicative systems, compare gē sēoð and hē sīeð. Identify each verb, class, tense, person and number. These are form-identification examples; add no complement.
% mdparagraph:P301:end

\end{minipage}\par\addvspace{4pt}

\clearpage
% mdparagraph:P302:begin
\section*{14. Sources and original reading notes}
% mdparagraph:P302:end

\begingroup\fontsize{10}{12.8}\selectfont
\setlength{\parskip}{1.5pt}
% mdparagraph:P303:begin
The core strong-verb systems follow Cook, \textbf{A First Book in Old English}, §§105–110, printed pp.63–67, with the earlier-class additions from §§101–104. Wright’s \textbf{Old English Grammar} and the checked dictionary entries corroborate the stated alternatives and lexical uses. Complete finite tables apply the ordinary endings to the checked selected stems; they do not claim a separate surviving quotation for every cell.
% mdparagraph:P303:end

% mdparagraph:P304:begin
The supplied participle \textbf{weaxen} is a rule-derived expansion of Cook’s and Wright’s stated Class VII model. It is not a whole-form attestation in the checked rows. Its status remains citation support only.
% mdparagraph:P304:end

% mdparagraph:P305:begin
The lexical constructions use Cook’s case rules and independently checked dictionary evidence for the ordinary verbs. The eight examples and the connected account are new compositions. Their food, companion, bird and fish references are ordinary applications of the stated uses. Diagnostic form practice supplies no additional complement syntax.
% mdparagraph:P305:end

% mdparagraph:P306:begin
The source readings preserve Cook’s Old English spelling, quantity and punctuation. Line breaks are joined and note cues are moved to this section. The first-day passage is an explicit revisit of the Unit 1 text; the second-day paragraph is the new continuation. All relevant original notes follow. References such as “53. 3.” point to Cook’s own grammar or dictionary, not to sections of this unit. The local help above supplies what is needed here.
% mdparagraph:P306:end

\reserve{7\baselineskip}
% mdparagraph:P307:begin
\subsection*{Cook’s notes to the first day}
% mdparagraph:P307:end

% source-note: FB123-n1
% mdparagraph:P308:begin
\textbf{P.123, note 1.} See gescieppan, and 18.
% mdparagraph:P308:end

% source-note: FB123-n2
% mdparagraph:P309:begin
\textbf{P.123, note 2.} The order is probably determined by the Latin: creavit Deus.
% mdparagraph:P309:end

% source-note: FB123-n3
% mdparagraph:P310:begin
\textbf{P.123, note 3.} 53. 3.
% mdparagraph:P310:end

% source-note: FB123-n4
% mdparagraph:P311:begin
\textbf{P.123, note 4.} See se.
% mdparagraph:P311:end

% source-note: FB123-n5
% mdparagraph:P312:begin
\textbf{P.123, note 5.} Lat. autem.
% mdparagraph:P312:end

% source-note: FB123-n6
% mdparagraph:P313:begin
\textbf{P.123, note 6.} See wesan.
% mdparagraph:P313:end

% source-note: FB123-n7
% mdparagraph:P314:begin
\textbf{P.123, note 7.} Plural, like Lat. tenebrae.
% mdparagraph:P314:end

% source-note: FB123-n8
% mdparagraph:P315:begin
\textbf{P.123, note 8.} Governs brādnesse.
% mdparagraph:P315:end

% source-note: FB123-n9
% mdparagraph:P316:begin
\textbf{P.123, note 9.} Genitive, dependent on brādnesse (153. i).
% mdparagraph:P316:end

% source-note: FB123-n10
% mdparagraph:P317:begin
\textbf{P.123, note 10.} See 166.
% mdparagraph:P317:end

% source-note: FB123-n11
% mdparagraph:P318:begin
\textbf{P.123, note 11.} wæs gefęred = Lat. ferebatur. See gefęrian.
% mdparagraph:P318:end

% source-note: FB123-n12
% mdparagraph:P319:begin
\textbf{P.123, note 12.} See wæter, and 47. 1, 6.
% mdparagraph:P319:end

% source-note: FB123-n13
% mdparagraph:P320:begin
\textbf{P.123, note 13.} See cweðan.
% mdparagraph:P320:end

% source-note: FB123-n14
% mdparagraph:P321:begin
\textbf{P.123, note 14.} See geweorðan, and 193. a.
% mdparagraph:P321:end

% source-note: FB123-n15
% mdparagraph:P322:begin
\textbf{P.123, note 15.} See weorðan.
% mdparagraph:P322:end

% source-note: FB123-n16
% mdparagraph:P323:begin
\textbf{P.123, note 16.} Wearð geworht = facta est. See gewyrcean.
% mdparagraph:P323:end

% source-note: FB123-n17
% mdparagraph:P324:begin
\textbf{P.123, note 17.} See geseon.
% mdparagraph:P324:end

% source-note: FB123-n18
% mdparagraph:P325:begin
\textbf{P.123, note 18.} See hē.
% mdparagraph:P325:end

% source-note: FB124-n1
% mdparagraph:P326:begin
\textbf{P.124, note 1.} See wesan.
% mdparagraph:P326:end

% source-note: FB124-n2
% mdparagraph:P327:begin
\textbf{P.124, note 2.} See gedǣlan.
% mdparagraph:P327:end

% source-note: FB124-n3
% mdparagraph:P328:begin
\textbf{P.124, note 3.} See se.
% mdparagraph:P328:end

% source-note: FB124-n4
% mdparagraph:P329:begin
\textbf{P.124, note 4.} See p.123, note 7.
% mdparagraph:P329:end

% source-note: FB124-n5
% mdparagraph:P330:begin
\textbf{P.124, note 5.} See hātan, and 189, note.
% mdparagraph:P330:end

% source-note: FB124-n6
% mdparagraph:P331:begin
\textbf{P.124, note 6.} Wæs geworden = factum est. See geweorðan.
% mdparagraph:P331:end

% source-note: FB124-n7
% mdparagraph:P332:begin
\textbf{P.124, note 7.} Lat. dies unus.
% mdparagraph:P332:end

\reserve{7\baselineskip}
% mdparagraph:P333:begin
\subsection*{Cook’s notes to the second day}
% mdparagraph:P333:end

% source-note: FB124-n8
% mdparagraph:P334:begin
\textbf{P.124, note 8.} See cweðan.
% mdparagraph:P334:end

% source-note: FB124-n9
% mdparagraph:P335:begin
\textbf{P.124, note 9.} Lat. quoque.
% mdparagraph:P335:end

% source-note: FB124-n10
% mdparagraph:P336:begin
\textbf{P.124, note 10.} See geweorðan, and 193. a.
% mdparagraph:P336:end

% source-note: FB124-n11
% mdparagraph:P337:begin
\textbf{P.124, note 11.} See wæter, and 47. 1, 6.
% mdparagraph:P337:end

% source-note: FB124-n12
% mdparagraph:P338:begin
\textbf{P.124, note 12.} See tōtwǣman.
% mdparagraph:P338:end

% source-note: FB124-n13
% mdparagraph:P339:begin
\textbf{P.124, note 13.} See 87. d.
% mdparagraph:P339:end

% source-note: FB124-n14
% mdparagraph:P340:begin
\textbf{P.124, note 14.} Past part. of gedōn.
% mdparagraph:P340:end

% source-note: FB124-n15
% mdparagraph:P341:begin
\textbf{P.124, note 15.} See 173.
% mdparagraph:P341:end

% source-note: FB124-n16
% mdparagraph:P342:begin
\textbf{P.124, note 16.} Lat. secundus.
% mdparagraph:P342:end

\endgroup
\end{document}
